% DERIVED from formal_layer.json — read-only, do not edit.
\begin{assumptionv}{ass:uniform}[Uniform covariate distribution]
Let \(\lambda\) denote Lebesgue probability measure on \([0,1]\). The covariate marginal \(P_X\) satisfies
\[
P_X=\lambda.
\]
\end{assumptionv}

\begin{assumptionv}{ass:overlap}[Fixed propensity overlap]
The continuous propensity representative \(e_P\) satisfies
\[
e_P([0,1])\subseteq[1/4,3/4].
\]
\end{assumptionv}

\begin{assumptionv}{ass:propensity-smoothness}[Propensity Hölder smoothness]
Let \(\mathcal H^\alpha(20)\) denote the continuous functions on \([0,1]\) with supremum norm and \(\alpha\)-Hölder constant at most \(20\), and let \([f]_\lambda\) denote the equivalence class of \(f\) under almost-everywhere equality for Lebesgue probability measure \(\lambda\). The propensity equivalence class \(\mathcal E_P\) satisfies
\[
\mathcal E_P\in\{[f]_\lambda:f\in\mathcal H^\alpha(20)\}.
\]
\end{assumptionv}

\begin{assumptionv}{ass:baseline-smoothness}[Baseline mean smoothness]
Let \(\mathcal H^\beta(20)\) denote the continuous functions on \([0,1]\) with supremum norm and \(\beta\)-Hölder constant at most \(20\), and let \([g]_\lambda\) denote the equivalence class of \(g\) under almost-everywhere equality for Lebesgue probability measure \(\lambda\). The baseline conditional arm-mean equivalence class \(\mathcal G_{0,P}\) satisfies
\[
\mathcal G_{0,P}\in\{[g]_\lambda:g\in\mathcal H^\beta(20)\}.
\]
\end{assumptionv}

\begin{assumptionv}{ass:effect-smoothness}[Mean effect smoothness]
Let \(\mathcal H^\gamma(20)\) denote the continuous functions on \([0,1]\) with supremum norm and \(\gamma\)-Hölder constant at most \(20\), and let \([f]_\lambda\) denote the equivalence class of \(f\) under almost-everywhere equality for Lebesgue probability measure \(\lambda\). With \(m_{0,P}\) the continuous baseline arm-mean representative, the treated conditional arm-mean equivalence class \(\mathcal G_{1,P}\) satisfies
\[
\mathcal G_{1,P}\in\{[m_{0,P}+t]_\lambda:t\in\mathcal H^\gamma(20)\}.
\]
\end{assumptionv}

\begin{assumptionv}{ass:baseline-cap}[Baseline mean bound]
The continuous baseline arm-mean representative \(m_{0,P}\) satisfies
\[
\|m_{0,P}\|_\infty\le1/2.
\]
\end{assumptionv}

\begin{assumptionv}{ass:effect-cap}[Uniform mean effect bound]
The mean effect \(\tau_P\), the canonical difference between the treated and baseline conditional means, satisfies
\[
\|\tau_P\|_\infty\le \frac12.
\]
\end{assumptionv}

\begin{assumptionv}{ass:raw-moment}[Uniform conditional moment bound]
Let \(\lambda\) denote Lebesgue probability measure on \([0,1]\), and let \(Q_{a,P}(dy\mid x)\) denote the conditional outcome probability kernel in treatment arm \(a\). The conditional raw moment of order \(p\) satisfies
\[
\max_{a\in\{0,1\}}
\operatorname*{ess\,sup}_{x\sim\lambda}
\int |y|^p\,Q_{a,P}(dy\mid x)\le 10.
\]
\end{assumptionv}

\begin{definitionv}{def:model}[Original-record model \(\mathcal M_v\)]
The model \(\mathcal M_v\) is the set of original-record probability laws \(P\) satisfying the following restrictions, indexed by \(v=(p,\alpha,\beta,\gamma)\in\mathcal V\), where \(\mathcal V\) denotes the admissible domain of public moment and primitive smoothness tuples:
\begin{itemize}
\item \textbf{(Covariate law.)} \cref{obj:ass:uniform}.
\item \textbf{(Overlap.)} \cref{obj:ass:overlap}.
\item \textbf{(Propensity smoothness.)} \cref{obj:ass:propensity-smoothness}.
\item \textbf{(Baseline smoothness.)} \cref{obj:ass:baseline-smoothness}.
\item \textbf{(Effect smoothness.)} \cref{obj:ass:effect-smoothness}.
\item \textbf{(Baseline bound.)} \cref{obj:ass:baseline-cap}.
\item \textbf{(Effect bound.)} \cref{obj:ass:effect-cap}.
\item \textbf{(Conditional moment bound.)} \cref{obj:ass:raw-moment}.
\end{itemize}
Here \(P_X\) denotes the covariate marginal of \(P\). For a function \(f\), write \([f]_\lambda\) for its equivalence class under Lebesgue almost-everywhere equality. Let \(\mathcal E_P\) denote the propensity equivalence class, and let \(\mathcal G_{0,P}\) and \(\mathcal G_{1,P}\) denote the baseline and treated conditional arm-mean equivalence classes. Their continuous representatives are defined by
\[
\mathcal E_P=[e_P]_\lambda,\qquad
\mathcal G_{0,P}=[m_{0,P}]_\lambda,\qquad
\mathcal G_{1,P}=[m_{1,P}]_\lambda,
\qquad
e_P,m_{0,P},m_{1,P}\in C([0,1]),
\]
where \(C([0,1])\) is the space of continuous real functions on \([0,1]\) with the uniform topology. The canonical mean effect is
\[
\tau_P=m_{1,P}-m_{0,P}.
\]
\end{definitionv}

\begin{assumptionv}{ass:null-constancy}[Constant mean effect restriction]
There exists \(c\in[-1/2,1/2]\) such that
\[
\tau_P(x)=c
\qquad\text{for every }x\in[0,1].
\]
\end{assumptionv}

\begin{definitionv}{def:null}[Constant-effect null $H_0(v)$]
\[
H_0(v)
=
\left\{
P\in\mathcal M_v:
\exists c\in[-1/2,1/2]\ 
\forall x\in[0,1],\ \tau_P(x)=c
\right\}.
\]
The model restrictions and null constancy condition are given in \cref{obj:ass:uniform,obj:ass:overlap,obj:ass:propensity-smoothness,obj:ass:baseline-smoothness,obj:ass:effect-smoothness,obj:ass:baseline-cap,obj:ass:effect-cap,obj:ass:raw-moment,obj:ass:null-constancy}.
\end{definitionv}

\begin{definitionv}{def:causal-completion}[Potential-outcome completion $\mathbb Q_P$]
The potential-outcome completion \(\mathbb Q_P\) is defined using Lebesgue probability measure \(\lambda\) on \([0,1]\) and the conditional outcome probability kernels \(Q_{a,P}(dy\mid x)\) in arms \(a\in\{0,1\}\) by
\[
\mathbb Q_P(dx,dy_0,dy_1,da)
=
\lambda(dx)\,Q_{0,P}(dy_0\mid x)\,Q_{1,P}(dy_1\mid x)\,
e_P(x)^a(1-e_P(x))^{1-a}.
\]
Writing \(Y(0)\) and \(Y(1)\) for the baseline and treated potential-outcome coordinates, respectively, the observed outcome satisfies
\[
Y=(1-A)Y(0)+AY(1).
\]
Changes to the conditional outcome kernels on \(\lambda\)-null sets leave \(\mathbb Q_P\) unchanged.
\end{definitionv}

\begin{definitionv}{def:original-record-experiment}[Original-record experiment \(\mathsf S_{n,P}\)]
The original-record experiment is the probability law
\[
\mathsf S_{n,P}=P^{\otimes n}\otimes\lambda
\]
on
\[
\bigl([0,1]\times\{0,1\}\times\mathbb R\bigr)^n\times[0,1],
\]
where \(P\) is the original-record law and \(\lambda\) is Lebesgue probability measure on \([0,1]\). Each original record has coordinates \(X\in[0,1]\), \(A\in\{0,1\}\), and \(Y\in\mathbb R\), representing the covariate, treatment, and observed outcome. Write
\[
\mathcal D_n=\bigl((X_i,A_i,Y_i)\bigr)_{i=1}^n
\]
for the independent sample and \(U\) for the independent public uniform randomization variable. For a randomized rejection-probability map \(\phi\), its rejection expectation is
\[
\mathsf R_n(P,\phi)
=
\int\phi(\mathcal D_n,u)\,
\mathsf S_{n,P}(d\mathcal D_n,du).
\]
\end{definitionv}

\begin{definitionv}{def:testing-risk}[Minimax testing risk $B_n(v,r)$]
\[
B_n(v,r)=
\inf_{\phi:\ \sup_{P\in H_0(v)}\mathsf R_n(P,\phi)\le 1/10}
\ \sup_{P\in\mathcal M_v:\ d(P)\ge r}
\bigl\{1-\mathsf R_n(P,\phi)\bigr\}.
\]
The parameters satisfy:
\begin{itemize}
\item \textbf{(Sample size.)} \(n\ge 2\).
\item \textbf{(Public tuple.)} \(v\in\mathcal V\), the admissible domain of moment and primitive smoothness tuples.
\item \textbf{(Separation.)} \(0<r<D_v\), where \(D_v\) is the supremum of \(d(P)\) over \(\mathcal M_v\).
\end{itemize}
Here \(\mathcal M_v\) is the original-record model, \(H_0(v)\) consists of its laws with an unknown constant mean effect, and \(d(P)\) is the population Lebesgue \(L_2\) distance of the effect from its average. The infimum ranges over randomized rejection-probability maps \(\phi\), and \(\mathsf R_n(P,\phi)\) is their expected rejection probability in the original-record experiment.
\end{definitionv}

\begin{definitionv}{def:critical-radius}[Capped critical radius $r_n^*(v)$]
\[
r_n^*(v)=
\inf\Bigl(
\bigl\{r\in(0,D_v):B_n(v,r)\le 1/10\bigr\}
\cup\{D_v\}
\Bigr).
\]
Here \(B_n(v,r)\) is the original-model minimax type-II error under the size constraint \(1/10\), and \(D_v\) is the supremum of the population effect distance over the original model.
\end{definitionv}

\begin{assumptionv}{ass:bounded-outcome}[Almost sure outcome bound]
The observed outcome \(Y\) satisfies
\[
P(|Y|\le 1)=1.
\]
\end{assumptionv}

\begin{definitionv}{def:bounded-model}[Bounded-outcome model $\mathcal M^{\mathrm b}_w$]
The bounded-outcome model \(\mathcal M^{\mathrm b}_w\) is defined for the matched smoothness tuple \(w=(\alpha,\beta,\gamma)\) in \(\mathcal W\), the admissible domain of matched smoothness tuples, by
\[
\mathcal M^{\mathrm b}_w
=
\left\{
P\in\mathcal M_{(2,\alpha,\beta,\gamma)}:
P(|Y|\le 1)=1
\right\},
\qquad
w=(\alpha,\beta,\gamma)\in\mathcal W.
\]
Its restrictions are given in \cref{obj:ass:uniform,obj:ass:overlap,obj:ass:propensity-smoothness,obj:ass:baseline-smoothness,obj:ass:effect-smoothness,obj:ass:baseline-cap,obj:ass:effect-cap,obj:ass:raw-moment,obj:ass:bounded-outcome}.
\end{definitionv}

\begin{assumptionv}{ass:binary-outcome}[Signed binary outcome support]
The observed outcome \(Y\) satisfies
\[
P\bigl(Y\in\{-1,1\}\bigr)=1.
\]
\end{assumptionv}

\begin{definitionv}{def:binary-model}[Signed-binary model $\mathcal M^{\mathrm{bin}}_w$]
The signed-binary model \(\mathcal M^{\mathrm{bin}}_w\) is defined for the matched smoothness tuple \(w=(\alpha,\beta,\gamma)\) in \(\mathcal W\), the admissible domain of matched smoothness tuples, by
\[
\mathcal M^{\mathrm{bin}}_w
=
\left\{
P\in\mathcal M_{(2,\alpha,\beta,\gamma)}:
P(Y\in\{-1,1\})=1
\right\},
\qquad
w=(\alpha,\beta,\gamma)\in\mathcal W.
\]
Its restrictions are given in \cref{obj:ass:uniform,obj:ass:overlap,obj:ass:propensity-smoothness,obj:ass:baseline-smoothness,obj:ass:effect-smoothness,obj:ass:baseline-cap,obj:ass:effect-cap,obj:ass:raw-moment,obj:ass:binary-outcome}.
\end{definitionv}

\begin{definitionv}{def:bounded-null}[Bounded constant-effect null $H_0^{\mathrm b}(w)$]
\[
H_0^{\mathrm b}(w)
=
\left\{
P\in\mathcal M^{\mathrm b}_w:
\exists c\in[-1/2,1/2]\ 
\forall x\in[0,1],\ \tau_P(x)=c
\right\}.
\]
The model restrictions and null constancy condition are given in \cref{obj:ass:uniform,obj:ass:overlap,obj:ass:propensity-smoothness,obj:ass:baseline-smoothness,obj:ass:effect-smoothness,obj:ass:baseline-cap,obj:ass:effect-cap,obj:ass:raw-moment,obj:ass:bounded-outcome,obj:ass:null-constancy}.
\end{definitionv}

\begin{definitionv}{def:binary-null}[Signed-binary constant-effect null $H_0^{\mathrm{bin}}(w)$]
\[
H_0^{\mathrm{bin}}(w)
=
\left\{
P\in\mathcal M^{\mathrm{bin}}_w:
\exists c\in[-1/2,1/2]\ 
\forall x\in[0,1],\ \tau_P(x)=c
\right\}.
\]
The model restrictions and null constancy condition are given in \cref{obj:ass:uniform,obj:ass:overlap,obj:ass:propensity-smoothness,obj:ass:baseline-smoothness,obj:ass:effect-smoothness,obj:ass:baseline-cap,obj:ass:effect-cap,obj:ass:raw-moment,obj:ass:binary-outcome,obj:ass:null-constancy}.
\end{definitionv}

\begin{definitionv}{def:bounded-risk}[Bounded-model testing risk]
The bounded-model minimax type-II error \(B_n^{\mathrm b}(w,r)\) is defined by
\[
B_n^{\mathrm b}(w,r)
=
\inf_{\phi:\,\sup_{P\in H_0^{\mathrm b}(w)}\mathsf R_n(P,\phi)\le 1/10}
\sup_{P\in\mathcal M^{\mathrm b}_w:\,d(P)\ge r}
\{1-\mathsf R_n(P,\phi)\},
\]
where the tests \(\phi\), rejection probabilities \(\mathsf R_n(P,\phi)\), and heterogeneity distances \(d(P)\) are as in \cref{obj:def:testing-risk}, and the bounded model and null are those of \cref{obj:def:bounded-model,obj:def:bounded-null}. The parameter domain is:
\begin{itemize}
\item \textbf{(Sample size.)} \(n\in\mathbb N\) with \(n\ge2\).
\item \textbf{(Smoothness.)} \(w=(\alpha,\beta,\gamma)\in\mathcal W\), where the admissible smoothness domain is
\[
\mathcal W
=
\{(\alpha,\beta,\gamma)\in\mathbb R^3:
0<\alpha\le1,\quad 0<\beta\le1,\quad 1/4\le\gamma\le1\}.
\]
\item \textbf{(Separation.)} \(0<r<D_w^{\mathrm b}\), where the bounded-model distance supremum is
\[
D_w^{\mathrm b}
=
\sup_{P\in\mathcal M^{\mathrm b}_w}d(P).
\]
\end{itemize}
\end{definitionv}

\begin{definitionv}{def:binary-risk}[Signed-binary minimax testing risk]
The signed-binary minimax type-II error \(B_n^{\mathrm{bin}}(w,r)\) is defined by
\[
B_n^{\mathrm{bin}}(w,r)
=
\inf_{\phi:\,\sup_{P\in H_0^{\mathrm{bin}}(w)}
\mathsf R_n(P,\phi)\le 1/10}
\sup_{P\in\mathcal M^{\mathrm{bin}}_w:\,d(P)\ge r}
\{1-\mathsf R_n(P,\phi)\},
\]
where the tests and their expected rejection probabilities \(\mathsf R_n(P,\phi)\) follow \cref{obj:def:testing-risk}, and the signed-binary model and null are specified in \cref{obj:def:binary-model,obj:def:binary-null}. Here \(\lambda\) denotes Lebesgue probability measure on \([0,1]\), and the population heterogeneity distance is
\[
d(P)
=
\left(
\int_{[0,1]}
\left\{\tau_P(x)-\int_{[0,1]}\tau_P(z)\,d\lambda(z)\right\}^{2}
\,d\lambda(x)
\right)^{1/2}.
\]
The public parameter domain is:
\begin{itemize}
\item \textbf{(Sample size.)} \(n\in\mathbb N\) with \(n\ge2\).
\item \textbf{(Smoothness.)} \(w=(\alpha,\beta,\gamma)\in\mathcal W\), where
\[
\mathcal W=(0,1]\times(0,1]\times[1/4,1].
\]
\item \textbf{(Separation.)} \(0<r<D_w^{\mathrm{bin}}\), where the supremum of legal signed-binary heterogeneity distances is
\[
D_w^{\mathrm{bin}}
=
\sup_{P\in\mathcal M^{\mathrm{bin}}_w}d(P).
\]
\end{itemize}
\end{definitionv}

\begin{definitionv}{def:bounded-radius}[Bounded-model critical separation radius]
The capped bounded-model critical separation radius \(r_n^{*,\mathrm b}(w)\) is defined for \(n\in\mathbb N\) and \(w=(\alpha,\beta,\gamma)\in\mathbb R^3\). Write \(d(P)\) for the population Lebesgue \(L_2\) distance of the mean effect under \(P\) from its average, and define the bounded-model cap by
\[
D_w^{\mathrm b}
=
\sup_{P\in\mathcal M^{\mathrm b}_w} d(P),
\]
where \(\mathcal M^{\mathrm b}_w\) is the bounded model in \cref{obj:def:bounded-model}. With \(B_n^{\mathrm b}(w,r)\), the bounded-model testing risk defined in \cref{obj:def:bounded-risk}, set
\[
r_n^{*,\mathrm b}(w)
=
\inf\left(
\left\{r\in(0,D_w^{\mathrm b}):B_n^{\mathrm b}(w,r)\le \frac{1}{10}\right\}
\cup\{D_w^{\mathrm b}\}
\right).
\]
\end{definitionv}

\begin{definitionv}{def:binary-radius}[Capped signed-binary critical radius]
The capped signed-binary critical separation radius \(r_n^{*,\mathrm{bin}}(w)\), for sample size \(n\in\mathbb N\) and primitive smoothness tuple \(w=(\alpha,\beta,\gamma)\in\mathbb R^3\), is
\[
r_n^{*,\mathrm{bin}}(w)
=
\inf\left(
\left\{r\in(0,D_w^{\mathrm{bin}}):
B_n^{\mathrm{bin}}(w,r)\le \frac{1}{10}\right\}
\cup\{D_w^{\mathrm{bin}}\}
\right),
\]
where \(D_w^{\mathrm{bin}}\) is the supremum of the population heterogeneity distances over the signed-binary model in \cref{obj:def:binary-model}, and \(B_n^{\mathrm{bin}}(w,r)\) is the signed-binary testing risk in \cref{obj:def:binary-risk}.
\end{definitionv}

\begin{definitionv}{def:sharp-frontier-scales}[Separation scales \(\rho_n(v)\) and \(\rho_n^{\mathrm b}(w)\)]
The separation scales and lower multipliers are
\[
 \rho_n(v)=n^{-E(p)},\qquad
 \rho_n^{\mathrm b}(w)=n^{-E(2)},\qquad
 c_w^{\mathrm b}=c_{(2,w)},
\]
\[
 c_v=
 \begin{cases}
 c_0^{\mathrm e},&F(p)\ge1,\\
 \min\!\left\{k_v,\,
 c_0^{\mathrm e}N_{\mathrm{low}}(v)^{-(E_0(p)-E_4(p))}\right\},
 &F(p)<1.
 \end{cases}
\]
For an admissible public tuple \(v=(p,\alpha,\beta,\gamma)\in\mathcal V\) and its matched smoothness tuple \(w=(\alpha,\beta,\gamma)\), the quantities in these formulas are defined by
\[
 q=\frac{p-1}{p},\qquad S=\alpha+\beta,\qquad
 D_p=1+2\alpha+\frac{\beta}{q}+\frac{S}{2\gamma},
\]
\[
 E_0(p)=\frac{2\gamma q}{2\gamma+q},\qquad
 E_4(p)=\frac{2S}{D_p},\qquad
 F(p)=\frac{2\alpha}{p-1}+\frac{\beta}{q}+\frac{S}{2\gamma},
 \qquad E(p)=\min\{E_0(p),E_4(p)\},
\]
\[
 H_{\mathrm{fine}}=2^{40},\qquad
 N_{\mathrm{low}}(v)=
 \left\lceil\max\left\{2,32^{D_p\gamma/(2S)}\right\}\right\rceil,
\]
\[
 k_v=2^{-17}(2H_{\mathrm{fine}})^{-S},\qquad
 c_0^{\mathrm e}=\frac{4^{-\gamma}}{16\sqrt3}.
\]
The arguments \(p\) in \(E_0,E_4,F,E\) vary with \((\alpha,\beta,\gamma)\) fixed; \(c_{(2,w)}\) is \(c_v\) evaluated at \(v=(2,\alpha,\beta,\gamma)\).

The upper multipliers and sample-size thresholds are the fully evaluable attainment quantities:
\[
 C_v=C^{\mathrm{att}}_v,\qquad N_v=N^{\mathrm{att}}_v,\qquad
 C_w^{\mathrm b}=C^{\mathrm{att}}_{(2,w)},\qquad
 N_w^{\mathrm b}=N^{\mathrm{att}}_{(2,w)}.
\]
The test \(\phi^*_{n,v}\) is the total original-record test specified by the explicit score construction, and its bounded-model counterpart is
\[
 \phi^{*,\mathrm b}_{n,w}=\phi^*_{n,(2,w)}.
\]
\end{definitionv}

\begin{definitionv}{def:comparison-handle}[Capped-radius ratio $\Delta_n(v)$]
\[
\Delta_n(v)=\frac{r_n^*(v)}{r_n^{*,\mathrm b}(w)},
\qquad w=(\alpha,\beta,\gamma).
\]
The numerator and denominator are the capped critical separation radii for the original and bounded original-record experiments, respectively, at the same sample size \(n\) and matched smoothness tuple \(w\).
\end{definitionv}

\begin{definitionv}{def:tail-region}[Tail region \(\mathcal T\)]
The tail region \(\mathcal T\) consists of admissible public tuples \(v=(p,\alpha,\beta,\gamma)\), with moment exponent \(p\) and primitive smoothness exponents \(\alpha,\beta,\gamma\), for which
\[
\mathcal T
=
\left\{
v\in\mathcal V:
p<2,\ 
\lim_{n\to\infty}\Delta_n(v)=\infty
\right\},
\]
where \(\mathcal V\) is the admissible tuple domain and \(\Delta_n(v)\) is the original-to-bounded capped-radius ratio at matched smoothness.
\end{definitionv}

\begin{definitionv}{def:equality-region}[Equality region \(\mathcal E\)]
The equality region \(\mathcal E\) is defined using \(\mathcal V\), the admissible domain of public moment and smoothness tuples, by
\[
\mathcal E
=
\left\{
v\in\mathcal V:
\sup_{n\ge 2}\Delta_n(v)<\infty
\right\},
\]
where \(\Delta_n(v)\) is the original-to-bounded capped-radius ratio at matched smoothness.
\end{definitionv}

\begin{propositionv}{prop:bounded-submodel-comparison}[Bounded and binary submodel comparison]
Let \(v=(p,\alpha,\beta,\gamma)\) be an admissible parameter tuple as in \cref{obj:def:model}, and let \(w=(\alpha,\beta,\gamma)\). Use the model classes in \cref{obj:def:binary-model,obj:def:bounded-model,obj:def:model}. Let \(\lambda\) denote Lebesgue probability measure on \([0,1]\), and define the heterogeneity distance, its model suprema, and the witness distance by
\[
d(P)=
\left[
\int_{[0,1]}
\left(\tau_P(x)-\int_{[0,1]}\tau_P(z)\,\lambda(dz)\right)^2
\,\lambda(dx)
\right]^{1/2},
\]
\[
D_w^{\mathrm{bin}}=\sup_{P\in\mathcal M^{\mathrm{bin}}_w}d(P),
\qquad
D_w^{\mathrm b}=\sup_{P\in\mathcal M^{\mathrm b}_w}d(P),
\qquad
D_v=\sup_{P\in\mathcal M_v}d(P),
\qquad
d_0=\frac{1}{8\sqrt{2}}.
\]
Then
\[
\mathcal M^{\mathrm{bin}}_w\subseteq\mathcal M^{\mathrm b}_w
\subseteq\mathcal M_v,
\qquad
D_w^{\mathrm{bin}}=D_w^{\mathrm b}=D_v\ge d_0.
\]
For every integer \(n\ge2\) and every real \(r\) satisfying \(0<r<D_v\), the testing risks of \cref{obj:def:binary-risk,obj:def:bounded-risk,obj:def:testing-risk} satisfy
\[
B_n^{\mathrm{bin}}(w,r)=B_n^{\mathrm b}(w,r)\le B_n(v,r).
\]
For every integer \(n\ge2\), the capped critical radii of \cref{obj:def:binary-radius,obj:def:bounded-radius,obj:def:critical-radius} satisfy
\[
0<r_n^{*,\mathrm{bin}}(w)
=r_n^{*,\mathrm b}(w)\le r_n^*(v).
\]

For every \(P\in\mathcal M_{(2,\alpha,\beta,\gamma)}\), its signed-binary conversion belongs to \(\mathcal M^{\mathrm{bin}}_w\). This conversion retains the distribution of \((X,A)\) and uses conditional signed-binary outcomes whose probabilities of \(+1\) in arms \(0\) and \(1\), respectively, are
\[
\frac{1+m_{0,P}(x)}{2},
\qquad
\frac{1+m_{1,P}(x)}{2},
\]
with the complementary probabilities assigned to \(-1\). It preserves the propensity \(e_P\), the baseline mean \(m_{0,P}\), the effect \(\tau_P\), and the heterogeneity distance \(d(P)\). The converted law satisfies the null constancy condition of \cref{obj:ass:null-constancy} if and only if the original law does.

For every \(n\in\mathbb N\), every test \(\phi\) in the sense of \cref{obj:def:testing-risk}, and every \(P\in\mathcal M^{\mathrm b}_w\), independently convert each original record by retaining \((X,A)\) and drawing an outcome sign with conditional probabilities
\[
\Pr(+1\mid X,A,Y)=\frac{1+Y}{2},
\qquad
\Pr(-1\mid X,A,Y)=\frac{1-Y}{2}.
\]
Integrating this conversion randomization into \(\phi\) gives a rejection-probability map on the original experiment. Its expected rejection probability under \(P\) equals the expected rejection probability of \(\phi\) under the signed-binary conversion of \(P\).
\end{propositionv}

\begin{definitionv}{def:projection}[Histogram projection \(\Pi_J\)]
For a positive integer \(J\), define the histogram cells, feature vector, and projection kernel by
\[
I_{j,J}=
\begin{cases}
[(j-1)/J,j/J),&1\le j<J,\\
[(J-1)/J,1],&j=J,
\end{cases}
\qquad
F_J(x)=\bigl(\sqrt J\,\mathbf 1_{I_{j,J}}(x)\bigr)_{j=1}^J,
\]
\[
\Pi_J(x,z)=\langle F_J(x),F_J(z)\rangle.
\]
The corresponding projection of an integrable function \(f\) is
\[
(\Pi_J f)(x)=\int_0^1\Pi_J(x,z)f(z)\,dz.
\]
\end{definitionv}

\begin{definitionv}{def:blocks}[Evaluation blocks \(\mathcal I_1,\mathcal I_2\)]
The two evaluation blocks and their common size are
\[
s_n=\lfloor n/2\rfloor,
\qquad
\mathcal I_1=\{1,\ldots,s_n\},
\qquad
\mathcal I_2=\{s_n+1,\ldots,2s_n\}.
\]
\end{definitionv}

\begin{definitionv}{def:score-family}[Score families \(H_{b,T}\) and \(U_{b,G,T}\)]
For \(n\ge4\), the single-record and ordered-pair score vectors are
\[
H_{b,T}(c)=
\frac{1}{s_n}\sum_{i\in\mathcal I_b}
F_{J_{\mathrm c}}(X_i)A_i\{\ell_T(Y_i)-c\},
\]
\[
U_{b,G,T}(c)=
\frac{1}{s_n(s_n-1)}
\sum_{\substack{i,j\in\mathcal I_b\\i\ne j}}
F_{J_{\mathrm c}}(X_i)A_iG(X_i,X_j)
\{\ell_T(Y_j)-cA_j\}.
\]
Here \(b\in\{1,2\}\), \(G:[0,1]^2\to\mathbb R\) is the pair-score kernel, and the clipping function is
\[
\ell_T(y)=\max\{-T,\min\{y,T\}\}.
\]
For \(n=2,3\), set
\[
H_{b,T}(c)=U_{b,G,T}(c)=0.
\]
\end{definitionv}

\begin{algorithmv}{def:explicit-score-ledger}[Explicit score and ledgers \(Z_b(c),\mathfrak b_{n,v},\Lambda_{n,v}\)]
Inputs are the original records and the public tuple \(v=(p,\alpha,\beta,\gamma)\); \(H_{b,T}(c)\) and \(U_{b,G,T}(c)\) denote the blockwise single-record histogram score and ordered-pair correction score, respectively.
\begin{enumerate}
\item Define the public exponents
\[
q=\frac{p-1}{p},
\qquad
S=\alpha+\beta,
\qquad
t=\frac{2-p}{p-1},
\]
\[
E_0=\frac{2\gamma q}{2\gamma+q},
\qquad
E_4=\frac{2S}{1+2\alpha+\beta/q+S/(2\gamma)},
\qquad
E=\min\{E_0,E_4\}.
\]
For \(n=2,3\), set
\[
J_{\mathrm c}=J_{\mathrm f}=1,
\qquad
T_j=1\quad\text{for every cutoff},
\qquad
Z_b(c)=0,
\]
\[
\mathfrak b_{n,v}=\Lambda_{n,v}=h_{n,v}=0,
\qquad
\phi^*_{n,v}=0.
\]
The remaining steps apply for \(n\ge4\).
\item Let \(s_n\) be the evaluation-block size. Set
\[
s=s_n,
\qquad
a=s^{-E},
\qquad
M=J_{\mathrm c}
=
2^{\left\lceil\log_2(a^{-1/\gamma})\right\rceil},
\]
\[
K=J_{\mathrm f}
=
2^{\left\lceil\log_2\!\left(\max\{M,a^{-1/S}\}\right)\right\rceil}.
\]
Define upper dyadic rounding and the main clipping quantities by
\[
\mathcal R(x)=2^{\lceil\log_2x\rceil},
\quad x>0,
\qquad
T_0=\mathcal R(a^{-1/(p-1)}),
\qquad
V_0=32T_0^{2-p}.
\]
\item If \(S\ge\gamma\) or \(\alpha\ge q\), set
\[
Z_b(c)=H_{b,T_0}(c)-U_{b,\Pi_K,T_0}(c),
\]
\[
B=400K^{-S}+20T_0^{1-p},
\qquad
L_1=16\sqrt{V_0},
\qquad
L_2=4\sqrt{KV_0}.
\]
\item Otherwise, set
\[
\lambda_0=\frac{(p-1)^2}{4p},
\qquad
L=\log_2(K/M),
\qquad
R_j=2^jM\quad(0\le j\le L),
\]
\[
D_j=\Pi_{R_j}-\Pi_{R_{j-1}},
\qquad 1\le j\le L.
\]
For each increment define
\[
T_j=
\mathcal R\!\left(
\max\!\left\{
1,\,
\min\!\left\{
a^{-1/(p-1)},
\left[
\frac{R_j^{-\alpha}}{a(R_j/K)^{\lambda_0}}
\right]^{1/(p-1)}
\right\}
\right\}
\right),
\qquad
V_j=32T_j^{2-p},
\]
\[
a_j=40R_{j-1}^{-\alpha},
\qquad
d_j=
110R_{j-1}^{-\min\{\alpha,\beta,\gamma\}}
+20T_j^{1-p}.
\]
Here \(L\) is an integer, and empty sums equal zero. Set
\[
Z_b(c)=
H_{b,T_0}(c)-U_{b,\Pi_M,T_0}(c)
-\sum_{j=1}^{L}U_{b,D_j,T_j}(c),
\]
\[
B=
400K^{-S}+20T_0^{1-p}
+10\sum_{j=1}^{L}a_jT_j^{1-p},
\]
\[
L_1=
16\sqrt{V_0}
+8\sum_{j=1}^{L}\bigl(d_j+a_j\sqrt{V_j}\bigr),
\]
\[
L_2=
4\left(
\sqrt{MV_0}+\sum_{j=1}^{L}\sqrt{R_jV_j}
\right).
\]
\item In either branch, define the public bias ledger, covariance ledger, and rejection cutoff
\[
\mathfrak b_{n,v}=B,
\qquad
\Lambda_{n,v}
=
\frac{L_1^2}{s}
+\frac{2L_2^2}{s(s-1)},
\]
\[
h_{n,v}
=
\mathcal R\!\left(
2B^2+2^{16}\sqrt M\,\Lambda_{n,v}
\right).
\]
Use the public rank and outcome grids
\[
\left\{2^j:0\le j\le\left\lceil\log_2(2s^2)\right\rceil\right\},
\qquad
\left\{2^j:0\le j\le\left\lceil\log_2(2s)\right\rceil\right\},
\]
respectively. The cutoff is the upper endpoint of the two-point dyadic grid bracketing its unrounded formula.
\item Form the quadratic statistic
\[
W(c)=\langle Z_1(c),Z_2(c)\rangle.
\]
Evaluate its minimum over \([-1/2,1/2]\) at the two endpoints and, when the quadratic coefficient is positive and the vertex is interior, at the vertex. Output
\[
\phi^*_{n,v}
=
\mathbf 1\!\left\{
\min_{|c|\le1/2}W(c)>h_{n,v}
\right\}.
\]
Every sum is finite. Histogram and ordered-pair scores take their stipulated values at bin endpoints and are evaluated directly from the original records, without division by observed cell counts.
\end{enumerate}
\end{algorithmv}

\begin{theoremv}{thm:whole-null-calibration-resolved}[Whole-null calibration and power]
Let \(v=(p,\alpha,\beta,\gamma)\) be an admissible public tuple as in \cref{obj:def:model}, and let \(n\ge2\) be an integer. Use the ranks, cutoffs, scores, bias ledger \(\mathfrak b_{n,v}\), and covariance ledger \(\Lambda_{n,v}\) of \cref{obj:def:explicit-score-ledger}. Write \(\psi_{n,v}\) for the original-record test constructed there, and define its rejection expectation by
\[
\mathsf R_n(P,\psi_{n,v})
:=\mathbb E\!\left[\psi_{n,v}(\text{sample},\text{seed})\right],
\]
where the sample consists of \(n\) independent records with law \(P\), and the seed is independent and uniform on the unit interval. Define the population score mean by
\[
\theta_P(c):=\mathbb E_{P^{\otimes n}}[Z_b(c)].
\]
Let \(\lambda\) denote Lebesgue probability measure on \([0,1]\), and define the heterogeneity distance and its model supremum by
\[
d(P):=
\left\{\int_0^1
\left(\tau_P(x)-\int_0^1\tau_P(t)\,d\lambda(t)\right)^2
\,d\lambda(x)\right\}^{1/2},
\qquad
D_v:=\sup_{P\in\mathcal M_v}d(P).
\]

For every \(P\in\mathcal M_v\), each affine coefficient of each block score is square-integrable, and its scalar projection in any direction \(u\in\mathbb R^{J_{\mathrm c}}\) has variance at most
\[
\Lambda_{n,v}\|u\|_2^2.
\]
For each block, the absolute covariance between projections of its two affine coefficients in directions \(u,z\in\mathbb R^{J_{\mathrm c}}\) is at most
\[
\Lambda_{n,v}\|u\|_2\|z\|_2.
\]
The simultaneous profiling event for \(W(c)\) specified in \cref{obj:def:explicit-score-ledger}, with profiling multiplier \(C:=128\), is measurable and has probability at least \(0.9925\) under \(P^{\otimes n}\). If \(\Lambda_{n,v}=0\), then, almost surely, simultaneously for every \(|c|\le1/2\),
\[
W(c)=\|\theta_P(c)\|_2^2.
\]

For every \(P\in H_0(v)\), with the null class defined in \cref{obj:def:null},
\[
\inf_{|c|\le1/2}\|\theta_P(c)\|_2\le\mathfrak b_{n,v},
\qquad
\mathsf R_n(P,\psi_{n,v})\le0.0075.
\]
Define the calibrated separation by
\[
R_{\mathrm{att}}(n,v)
:=\frac{16}{3}
\left\{
16J_{\mathrm c}^{-\gamma}
+5\mathfrak b_{n,v}
+1024J_{\mathrm c}^{1/4}\sqrt{\Lambda_{n,v}}
\right\}.
\]
If \(n\ge4\) and \(R_{\mathrm{att}}(n,v)<D_v\), then every \(P\in\mathcal M_v\) satisfying \(d(P)\ge R_{\mathrm{att}}(n,v)\) obeys
\[
1-\mathsf R_n(P,\psi_{n,v})\le0.0075.
\]
For \(n=2\) or \(n=3\), \(\psi_{n,v}\) equals zero at every sample and seed.
\end{theoremv}

\begin{definitionv}{def:oracle-broad-model}[Continuous-carrier model \(\mathcal M_v^{\mathrm{e,br}}\)]
The class \(\mathcal M_v^{\mathrm{e,br}}\) consists of the original-record laws satisfying the six conditions below and admitting continuous representatives of their propensity and both conditional arm means. Here \(v=(p,\alpha,\beta,\gamma)\in\mathcal V\) is an admissible public tuple, \(\lambda\) is Lebesgue probability measure on \([0,1]\), and \(Q_{a,P}(dy\mid x)\) is the conditional outcome probability kernel in treatment arm \(a\).
\begin{itemize}
\item \textbf{(Uniform covariates.)} The covariate marginal satisfies
\[
 P_X=\lambda,
\]
as stipulated in \cref{obj:ass:uniform}.
\item \textbf{(Overlap.)} The continuous propensity representative satisfies
\[
 e_P([0,1])\subseteq[1/4,3/4],
\]
as stipulated in \cref{obj:ass:overlap}.
\item \textbf{(Effect smoothness.)} The continuous mean effect satisfies
\[
 \tau_P=m_{1,P}-m_{0,P}\in\mathcal H^\gamma(20),
\]
where \(\mathcal H^\gamma(20)\) consists of continuous functions on \([0,1]\) whose sup norm and \(\gamma\)-Hölder constant are at most \(20\), as stipulated in \cref{obj:ass:effect-smoothness}.
\item \textbf{(Baseline cap.)} The baseline mean satisfies
\[
 \|m_{0,P}\|_\infty\le1/2,
\]
as stipulated in \cref{obj:ass:baseline-cap}.
\item \textbf{(Effect cap.)} The mean effect satisfies
\[
 \|\tau_P\|_\infty\le1/2,
\]
as stipulated in \cref{obj:ass:effect-cap}.
\item \textbf{(Raw moment.)} The conditional outcome kernels satisfy
\[
 \max_{a\in\{0,1\}}\operatorname*{ess\,sup}_{x\sim\lambda}
 \int |y|^p Q_{a,P}(dy\mid x)\le10,
\]
as stipulated in \cref{obj:ass:raw-moment}.
\end{itemize}
For a function \(h\), write \([h]_\lambda\) for its equivalence class under Lebesgue almost-everywhere equality. Let \(\mathcal E_P\) denote the propensity equivalence class and let \(\mathcal G_{0,P},\mathcal G_{1,P}\) denote the conditional arm-mean equivalence classes. The continuous-representative requirement is
\[
 \exists f,g_0,g_1\in C([0,1]):\qquad
 \mathcal E_P=[f]_\lambda,\qquad
 \mathcal G_{0,P}=[g_0]_\lambda,\qquad
 \mathcal G_{1,P}=[g_1]_\lambda.
\]
Here \(C([0,1])\) is the space of continuous real functions on the unit interval with the uniform topology. These representatives are the existing functionals \(e_P,m_{0,P},m_{1,P}\). Membership depends on \(p\) and \(\gamma\), with the full tuple \(v\) retained as the public index.

Writing \(d(P)\) for the population Lebesgue \(L_2\) distance of \(\tau_P\) from its average, define
\[
 D_v^{\mathrm{e,br}}=\sup_{P\in\mathcal M_v^{\mathrm{e,br}}}d(P).
\]
\end{definitionv}

\begin{definitionv}{def:oracle-broad-null}[Constant-effect null \(H_0^{\mathrm{e,br}}(v)\)]
The constant-effect null is
\[
 H_0^{\mathrm{e,br}}(v)=
 \left\{P\in\mathcal M_v^{\mathrm{e,br}}:
 \exists c\in[-1/2,1/2]\ \forall x\in[0,1],\
 \tau_P(x)=c\right\}.
\]
Membership in the continuous-carrier model imposes the restrictions in
\cref{obj:ass:uniform,obj:ass:overlap,obj:ass:effect-smoothness,obj:ass:baseline-cap,obj:ass:effect-cap,obj:ass:raw-moment}.
The unknown-constant effect condition is that of \cref{obj:ass:null-constancy}.
\end{definitionv}

\begin{definitionv}{def:oracle-broad-risk}[Supplied-propensity risk \(B_n^{\mathrm{e,br}}(v,r)\)]
The supplied-propensity minimax type-II error is
\[
 B_n^{\mathrm{e,br}}(v,r)=
 \inf_{\phi^{\mathrm e}:\,
     \sup_{P\in H_0^{\mathrm{e,br}}(v)}
       \mathsf R_n^{\mathrm e}(P,\phi^{\mathrm e})\le1/10}
 \ \sup_{P\in\mathcal M_v^{\mathrm{e,br}}:\,d(P)\ge r}
       \left\{1-\mathsf R_n^{\mathrm e}(P,\phi^{\mathrm e})\right\},
\]
for parameters satisfying
\begin{itemize}
\item \textbf{(Sample size.)} \(n\ge2\).
\item \textbf{(Public tuple.)} \(v\in\mathcal V\), the admissible domain of moment and primitive smoothness tuples.
\item \textbf{(Separation.)} \(0<r<D_v^{\mathrm{e,br}}\), where \(D_v^{\mathrm{e,br}}\) is the supremum of the population effect distance over \(\mathcal M_v^{\mathrm{e,br}}\).
\end{itemize}
Here \(d(P)\) is the population Lebesgue \(L_2\) distance of the mean effect from its average. The admissible tests \(\phi^{\mathrm e}\) are precisely the jointly Borel maps taking a supplied continuous function, an original-record sample, and a public randomization variable to a rejection probability in \([0,1]\). Their rejection expectation is
\[
 \mathsf R_n^{\mathrm e}(P,\phi^{\mathrm e})
 =\int \phi^{\mathrm e}(e_P,\mathcal D_n,u)\,
       P^{\otimes n}(d\mathcal D_n)\,du.
\]
The sample \(\mathcal D_n\) consists of \(n\) independent original records with law \(P\); the integration in \(u\) is over an independent uniform variable on \([0,1]\). The supplied function \(e_P\) is the whole true continuous propensity.
\end{definitionv}

\begin{definitionv}{def:oracle-broad-radius}[Capped radius \(r_n^{*,\mathrm{e,br}}(v)\)]
The capped supplied-propensity critical radius is
\[
 r_n^{*,\mathrm{e,br}}(v)=
 \inf\left(
 \left\{r\in(0,D_v^{\mathrm{e,br}}):
 B_n^{\mathrm{e,br}}(v,r)\le1/10\right\}
 \cup\left\{D_v^{\mathrm{e,br}}\right\}\right).
\]
Here \(D_v^{\mathrm{e,br}}\) is the supremum of the population effect distance over the continuous-carrier model and serves as the saturation value in the capped-radius convention.
\end{definitionv}

\begin{algorithmv}{def:oracle-handle}[Supplied-propensity test \(\phi^{*,\mathrm e}_{n,v}\)]
Inputs are \(n\ge2\) original records, the public tuple \(v=(p,\alpha,\beta,\gamma)\), and a supplied continuous propensity \(e\).
\begin{enumerate}
\item Set the block size, moment exponent, separation exponent, and target scale to
\[
s=\lfloor n/2\rfloor,
\qquad
q=\frac{p-1}{p},
\qquad
E_0=\frac{2\gamma q}{2\gamma+q},
\qquad
a=s^{-E_0}.
\]
Use the evaluation blocks
\[
\mathcal I_1=\{1,\ldots,s\},
\qquad
\mathcal I_2=\{s+1,\ldots,2s\}.
\]
\item Select the clipping cutoff and histogram rank
\[
T=2^{\left\lceil\log_2(a^{-1/(p-1)})\right\rceil},
\qquad
J=2^{\left\lceil\log_2(a^{-1/\gamma})\right\rceil}.
\]
Define outcome clipping and the overlap-clipped propensity by
\[
\ell_T(y)=\max\{-T,\min\{y,T\}\},
\qquad
\widetilde e(x)=\max\{1/4,\min\{e(x),3/4\}\}.
\]
For every legal supplied true propensity, \(\widetilde e=e\).
\item Form the original-record inverse-propensity scores
\[
R_i=
\frac{A_i\ell_T(Y_i)}{\widetilde e(X_i)}
-
\frac{(1-A_i)\ell_T(Y_i)}{1-\widetilde e(X_i)}.
\]
Let \(F_J\) be the normalized histogram feature vector. Its constant direction and the centered block vectors are
\[
u=J^{-1/2}(1,\ldots,1)^T,
\qquad
V_b=\frac1s\sum_{i\in\mathcal I_b}
\bigl(F_J(X_i)-u\bigr)R_i,
\quad b\in\{1,2\}.
\]
These vectors represent the block scores obtained by applying the constant-centered histogram projection \(\Pi_J-1\).
\item Set the public clipping-bias and covariance bounds to
\[
b=20T^{1-p},
\qquad
\Lambda=\frac{80T^{2-p}}s.
\]
Output the rejection rule
\[
\phi^{*,\mathrm e}_{n,v}
=
\mathbf 1\!\left\{
\langle V_1,V_2\rangle
>
2b^2+1024\sqrt J\,\Lambda
\right\}.
\]
Its calibration and necessity properties are given by its singleton and canonical quadratic projections and the normalized common-propensity paired-tent converse.
\end{enumerate}
\end{algorithmv}

\begin{definitionv}{def:oracle-risk}[Supplied-propensity risk $B_n^{\mathrm e}(v,r)$]
\[
B_n^{\mathrm e}(v,r)=
\inf_{\phi^{\mathrm e}:\ \sup_{P\in H_0(v)}
\mathsf R_n^{\mathrm e}(P,\phi^{\mathrm e})\le 1/10}
\ \sup_{P\in\mathcal M_v:\ d(P)\ge r}
\bigl\{1-\mathsf R_n^{\mathrm e}(P,\phi^{\mathrm e})\bigr\}.
\]
The parameters satisfy:
\begin{itemize}
\item \textbf{(Sample size.)} \(n\ge 2\).
\item \textbf{(Public tuple.)} \(v\in\mathcal V\), the admissible domain of moment and primitive smoothness tuples.
\item \textbf{(Separation.)} \(0<r<D_v\), where \(D_v\) is the supremum of \(d(P)\) over \(\mathcal M_v\).
\end{itemize}
Here \(\mathcal M_v\) is the original-record model, \(H_0(v)\) consists of its laws with an unknown constant mean effect, and \(d(P)\) is the population Lebesgue \(L_2\) distance of the effect from its average. The infimum ranges over randomized rejection-probability maps \(\phi^{\mathrm e}\) receiving the whole true continuous propensity, and \(\mathsf R_n^{\mathrm e}(P,\phi^{\mathrm e})\) is their expected rejection probability with that propensity supplied.
\end{definitionv}

\begin{definitionv}{def:oracle-radius}[Supplied-propensity radius $r_n^{*,\mathrm e}(v)$]
\[
r_n^{*,\mathrm e}(v)=
\inf\Bigl(
\bigl\{r\in(0,D_v):B_n^{\mathrm e}(v,r)\le 1/10\bigr\}
\cup\{D_v\}
\Bigr).
\]
Here \(B_n^{\mathrm e}(v,r)\) is the supplied-propensity minimax type-II error on the original model under the size constraint \(1/10\), and \(D_v\) is the supremum of the population effect distance over that model.
\end{definitionv}

\begin{definitionv}{def:preservation-region}[Propensity preservation region $\mathcal R$]
\[
\mathcal R=
\left\{
v\in\mathcal V:
\sup_{n\ge 2}
\frac{r_n^*(v)}{r_n^{*,\mathrm e}(v)}
<\infty
\right\}.
\]
Here \(\mathcal V\) is the admissible domain of public moment and primitive smoothness tuples. The radii \(r_n^*(v)\) and \(r_n^{*,\mathrm e}(v)\) are the capped critical separation radii on the same original-record model with unknown and supplied true continuous propensity, respectively.
\end{definitionv}

\begin{lemmav}{lem:causal-nonempty}[Nonemptiness and causal identification]
Fix an admissible parameter tuple \(v=(p,\alpha,\beta,\gamma)\) as in \cref{obj:def:model}. Let \(\lambda\) denote Lebesgue probability measure on \([0,1]\), and define the heterogeneity distance, its model supremum, and the witness distance by
\[
d(P)=
\left\{\int_0^1
\left(\tau_P(x)-\int_0^1\tau_P(z)\,d\lambda(z)\right)^2
\,d\lambda(x)\right\}^{1/2},
\qquad
D_v=\sup_{P\in\mathcal M_v}d(P),
\qquad
d_0=\frac{1}{8\sqrt{2}}.
\]
The set \(\{d(P):P\in\mathcal M_v\}\) is bounded above, and
\[
d_0\le D_v\le\frac12.
\]

For every \(P\in\mathcal M_v\), with model membership understood as in \cref{obj:def:model}, the continuous propensity, baseline mean, and contrast representatives are determined uniquely by the observed law: every other observational representation of the same probability law has exactly the same \(e_P\), \(m_{0,P}\), and \(\tau_P\), pointwise on \([0,1]\). Consequently, the treated mean representative \(m_{1,P}=m_{0,P}+\tau_P\) is also uniquely determined, and
\[
|m_{1,P}(x)-m_{1,P}(z)|
\le 40|x-z|^{\min\{\beta,\gamma\}}
\qquad (x,z\in[0,1]).
\]

In the completion \(\mathbb Q_P\) of \cref{obj:def:causal-completion}, let \(Y(0)\) and \(Y(1)\) denote the baseline and treated potential-outcome coordinates. The distribution under \(\mathbb Q_P\) of
\[
\bigl(X,A,(1-A)Y(0)+AY(1)\bigr)
\]
is \(P\), and
\[
(Y(0),Y(1))\perp A\mid X.
\]
Both potential outcomes are integrable under \(\mathbb Q_P\), and
\[
\mathbb E_{\mathbb Q_P}\!\left[Y(1)-Y(0)\mid X\right]
=\tau_P(X)
\qquad \mathbb Q_P\text{-almost surely}.
\]
Moreover,
\[
d(P)=0\quad\Longleftrightarrow\quad P\in H_0(v),
\]
where \(H_0(v)\) is defined in \cref{obj:def:null}.

The law with uniform covariate distribution, propensity and mean functions
\[
e_P(x)=\frac12+\frac{\sin(2\pi x)}{10},
\qquad
m_{0,P}(x)=\frac{\cos(2\pi x)}8,
\qquad
\tau_P(x)=\frac{\cos(2\pi x)}8,
\]
and deterministic conditional outcome
\[
Y=m_{0,P}(X)+A\tau_P(X)
\]
belongs to \(\mathcal M_v\) and has \(d(P)=d_0\). Replacing its contrast function by \(\tau_P(x)=1/8\) for every \(x\in[0,1]\), while retaining the displayed propensity, baseline mean, and deterministic conditional-outcome construction, yields a law in \(H_0(v)\).
\end{lemmav}

\begin{algorithmv}{def:calibration-handle}[Deterministic calibration handle]
The calibration handle, indexed by a sample size \(n\in\mathbb N\) and a public exponent tuple \(v=(p,\alpha,\beta,\gamma)\in\mathbb R^4\), is the ordered collection of the following seven components. Its branch, ranks, cutoffs, ledgers, and public grids are the deterministic selections from \(n\) and \(v\) in \cref{obj:def:explicit-score-ledger}.

\begin{enumerate}
\item The tuning collection
\[
\bigl(M,K,T_0,(T_j)_{j=1}^{L}\bigr),
\qquad M=J_{\mathrm c},\quad K=J_{\mathrm f},
\]
with the number of increments \(L\) specified in \cref{obj:def:explicit-score-ledger}.

\item The block-score family, taking a block index, a real candidate constant \(c\), and a dataset of \(n\) original records to a vector in \(\mathbb R^M\). For \(n\ge4\), it is
\[
Z_b(c)=
\begin{cases}
H_{b,T_0}(c)-U_{b,\Pi_K,T_0}(c),
&\text{single-correction branch},\\[1mm]
H_{b,T_0}(c)-U_{b,\Pi_M,T_0}(c)
-\displaystyle\sum_{j=1}^{L}U_{b,D_j,T_j}(c),
&\text{multiresolution branch},
\end{cases}
\]
where the block and kernel conventions are those of
\cref{obj:def:blocks,obj:def:score-family,obj:def:explicit-score-ledger}.
For \(n<4\), \(Z_b(c)=0\).

\item The two affine-coefficient maps, indexed by the block and a binary coefficient selector, with values
\[
Z_{b,0}=Z_b(0),
\qquad
Z_{b,A}=Z_b(0)-Z_b(1).
\]
Thus, for every real \(c\) and every dataset,
\[
Z_b(c)=Z_{b,0}-cZ_{b,A}.
\]

\item The deterministic ledger triple
\[
\bigl(\mathfrak b_{n,v},\Lambda_{n,v},h_{n,v}\bigr),
\]
whose bias, covariance, and rejection-cutoff selections are exactly those of
\cref{obj:def:explicit-score-ledger}.

\item The dataset-to-real profiling map
\[
\inf_{c\in[-1/2,1/2]}W(c),
\qquad
W(c)=\langle Z_1(c),Z_2(c)\rangle,
\]
using the two-block indexing convention of
\cref{obj:def:explicit-score-ledger}.

\item The original-record rejection rule specified in
\cref{obj:def:explicit-score-ledger}, namely, for \(n\ge4\),
\[
\mathbf 1\!\left\{
\inf_{c\in[-1/2,1/2]}W(c)>h_{n,v}
\right\},
\]
and the zero rejection rule for \(n<4\).

\item The ordered triple of finite public rank, outcome-cutoff, and rejection-cutoff grids selected in \cref{obj:def:explicit-score-ledger}. The first two grids depend on \(n\), and the third depends on \(n\) and \(v\).
\end{enumerate}
\end{algorithmv}

\begin{algorithmv}{def:mean-handle}[Original score mean handle]
The original score mean handle is the ordered pair of multiresolution and single-correction handles associated with the two score members of \cref{obj:def:calibration-handle}. Its parameters are an original-record law \(P\), ranks \(J,K,L\in\mathbb N\), a main cutoff \(T_{\mathrm{main}}\in\mathbb R\), and increment cutoffs \(T_j^{\mathrm{inc}}\in\mathbb R\) indexed by \(1\le j\le L\). Each branch retains, in order, a mean map indexed by \(c\in\mathbb R\), coefficient means, singleton projections, canonical pair projections, four population contractions for each coefficient pair, and a sample statistic indexed by \(n\) and \(c\).

The multiresolution branch uses coarse rank \(J_{\mathrm c}=J\) and fine rank \(J_{\mathrm f}=2^LJ\); the single-correction branch uses
\[
J_{\mathrm c}=J,\qquad J_{\mathrm f}=K.
\]
Both branches use \(T_{\mathrm{main}}\); the multiresolution branch also uses the increment cutoffs \(T_j^{\mathrm{inc}}\), \(1\le j\le L\). Define the projected propensity and propensity weight by
\[
e_K=\Pi_{J_{\mathrm f}}e_P,
\qquad
w_K(x)=e_P(x)\{1-e_K(x)\}.
\]
The retained mean map is
\[
\int_0^1 F_{J_{\mathrm c}}(x)
\left[
w_K(x)\{\tau_P(x)-c\}
+\{e_P(x)-e_K(x)\}m_{0,P}(x)
\right]\,dx
\]
plus the exact multiresolution or single-correction clipping-tail remainder, respectively, specified in \cref{obj:lem:original-score-mean-ledger}.

This retained map equals the population expectation of the corresponding corrected block score for every \(c\in\mathbb R\) and either evaluation block under the following conditions:
\begin{itemize}
\item \textbf{(Model.)} The public tuple \(v\) is admissible and \(P\in\mathcal M_v\), as specified in \cref{obj:def:model}.
\item \textbf{(Sample and coarse rank.)} \(n\ge4\) and \(J>0\).
\item \textbf{(Main cutoff.)} \(T_{\mathrm{main}}\ge1\).
\item \textbf{(Multiresolution cutoffs.)} In the multiresolution branch, \(T_j^{\mathrm{inc}}\ge1\) for every integer \(1\le j\le L\).
\item \textbf{(Single-correction refinement.)} In the single-correction branch, \(K>0\) and \(J\mid K\).
\end{itemize}

For each affine coefficient \(r\in\{0,1\}\), let \(h_r\) and \(g_r\) be the selected branch's raw single-record and symmetric pair kernels from \cref{obj:def:calibration-handle}. All expectations below use independent original records with law \(P\). Define the centered pair-kernel singleton projection by
\[
g_{r,1}(o)
=
\mathbb E[g_r(o,O')]
-\mathbb E[g_r(O,O')].
\]
The retained coefficient means, singleton projections, and canonical pair projections are, respectively,
\[
\mathbb E[h_r(O)]+\mathbb E[g_r(O,O')],
\]
\[
h_r(o)-\mathbb E[h_r(O)]+2g_{r,1}(o),
\]
\[
g_r(o,z)-\mathbb E[g_r(O,O')]
-g_{r,1}(o)-g_{r,1}(z).
\]
For every \(r,r'\in\{0,1\}\), the retained ordered quadruple of population contractions is defined by
\[
h_r\!\times h_{r'}
=
\mathbb E\langle h_r(O_1),h_{r'}(O_2)\rangle,
\]
\[
h_r\!\times g_{r'}
=
\mathbb E\langle h_r(O_1),g_{r'}(O_2,O_3)\rangle,
\]
\[
g_r\!\times h_{r'}
=
\mathbb E\langle g_r(O_1,O_2),h_{r'}(O_3)\rangle,
\]
\[
g_r\!\times g_{r'}
=
\mathbb E\langle g_r(O_1,O_2),g_{r'}(O_3,O_4)\rangle,
\]
where \(O_1,O_2,O_3,O_4\) are independent with law \(P\).

For each candidate constant \(c\in\mathbb R\), set
\[
h_c=h_0-ch_1,
\qquad
g_c=g_0-cg_1.
\]
For \(n\ge4\), use the block size \(s_n=\lfloor n/2\rfloor\) and the two deterministic blocks \(\mathcal I_1,\mathcal I_2\) of \cref{obj:def:blocks}. Define
\[
Z_b(c)=
\frac1{s_n}\sum_{i\in\mathcal I_b}h_c(O_i)
+
\frac1{s_n(s_n-1)}
\sum_{\substack{i,j\in\mathcal I_b\\i\ne j}}
g_c(O_i,O_j),
\qquad b\in\{1,2\}.
\]
The retained sample output is the single function \(c\mapsto W(c)\), given by
\[
\begin{aligned}
W(c)
&=\langle Z_1(c),Z_2(c)\rangle\\
&=\frac1{s_n^2}
\sum_{\substack{i\in\mathcal I_1\\j\in\mathcal I_2}}
\langle h_c(O_i),h_c(O_j)\rangle\\
&\quad+\frac1{s_n^2(s_n-1)}
\sum_{\substack{i\in\mathcal I_1\\j,k\in\mathcal I_2,\ j\ne k}}
\langle h_c(O_i),g_c(O_j,O_k)\rangle\\
&\quad+\frac1{s_n^2(s_n-1)}
\sum_{\substack{i,j\in\mathcal I_1,\ i\ne j\\k\in\mathcal I_2}}
\langle g_c(O_i,O_j),h_c(O_k)\rangle\\
&\quad+\frac1{s_n^2(s_n-1)^2}
\sum_{\substack{i,j\in\mathcal I_1,\ i\ne j\\k,l\in\mathcal I_2,\ k\ne l}}
\langle g_c(O_i,O_j),g_c(O_k,O_l)\rangle.
\end{aligned}
\]
For \(n<4\), its defining value is
\[
W(c)=0.
\]
The two mixed affine-coefficient contributions are retained through their sum in \(W(c)\).
\end{algorithmv}

\begin{lemmav}{lem:original-score-mean-ledger}[Original score mean bounds]
Let \(v=(p,\alpha,\beta,\gamma)\) be an admissible tuple and \(P\in\mathcal M_v\), as specified in \cref{obj:def:model}, and let \(n\ge4\) be an integer. Use the explicit score ledger's coarse rank \(M\), fine rank \(K\), corrected block scores \(Z_b(c)\), and deterministic bias bound \(B\).

Let \(\lambda\) denote Lebesgue probability measure on \([0,1]\), and define the projected propensity and its associated weight by
\[
e_K(x)=(\Pi_K e_P)(x),
\qquad
w_K(x)=e_P(x)\bigl(1-e_K(x)\bigr).
\]
Define the population heterogeneity distance by
\[
d(P)=
\left(
\int_0^1
\left[\tau_P(x)-\int_0^1\tau_P(z)\,d\lambda(z)\right]^2
\,d\lambda(x)
\right)^{1/2}.
\]

For every block index \(b\) and every \(c\in\mathbb R\), \(Z_b(c)\) is integrable under the original independent-record law \(P^{\otimes n}\). Its population mean is
\[
\theta_P(c)
=
\mathbb E_{P^{\otimes n}}\!\left[Z_b(c)\right].
\]
For every block index \(b\) and every \(c\in\mathbb R\) with \(|c|\le\tfrac12\),
\[
\left\|
\theta_P(c)
-
\int_0^1
w_K(x)\bigl(\tau_P(x)-c\bigr)F_M(x)\,d\lambda(x)
\right\|_2
\le B,
\qquad
\|\theta_P(c)\|_2
\ge
\frac{3}{16}d(P)-16M^{-\gamma}-B.
\]
Furthermore, if \(P\in H_0(v)\) as defined in \cref{obj:def:null}, then for every block index \(b\) and every \(c_0\in\mathbb R\) satisfying \(\tau_P(x)=c_0\) for all \(x\in[0,1]\),
\[
\|\theta_P(c_0)\|_2\le B.
\]
\end{lemmav}

\begin{lemmav}{lem:original-score-covariance-ledger}[Original score covariance bounds]
Let \(v=(p,\alpha,\beta,\gamma)\) be an admissible tuple of moment and smoothness exponents, and let \(P\in\mathcal M_v\), as specified in \cref{obj:def:model}. Let \(n\ge 2\) be an integer. Write the affine block score as
\[
Z_b(c)=Z_{b,0}-cZ_{b,A}.
\]
Under the independent-record sample law \(P^{\otimes n}\), for every evaluation block \(b\) and every \(u\in\mathbb R^M\), both \(u^T Z_{b,0}\) and \(u^T Z_{b,A}\) are square integrable and satisfy
\[
\operatorname{Var}_{P^{\otimes n}}(u^T Z_{b,0})
\le \Lambda_{n,v}\|u\|_2^2,
\qquad
\operatorname{Var}_{P^{\otimes n}}(u^T Z_{b,A})
\le \Lambda_{n,v}\|u\|_2^2.
\]
Moreover, for every evaluation block \(b\) and every \(u,z\in\mathbb R^M\),
\[
\left|
\operatorname{Cov}_{P^{\otimes n}}
\bigl(u^T Z_{b,0},z^T Z_{b,A}\bigr)
\right|
\le \Lambda_{n,v}\|u\|_2\|z\|_2.
\]
\end{lemmav}

\begin{lemmav}{lem:uniform-affine-profile-bound}[Uniform affine profile bound]
Let \(v=(p,\alpha,\beta,\gamma)\) be an admissible public tuple as in \cref{obj:def:model}, and let \(n\ge 2\) be an integer. For every \(P\in\mathcal M_v\), consider the explicit affine score construction, with block scores \(Z_1(c)\) and \(Z_2(c)\), coarse histogram rank \(J_{\mathrm c}\), and public covariance ledger \(\Lambda_{n,v}\). Define the score inner product and its population target by
\[
W(c)=\langle Z_1(c),Z_2(c)\rangle,
\qquad
\theta_P(c)=\mathbb E_{P^{\otimes n}}[Z_1(c)].
\]
The event on which, simultaneously for every \(c\in[-1/2,1/2]\),
\[
\bigl|W(c)-\|\theta_P(c)\|_2^2\bigr|
\le
128\left\{
\sqrt{\Lambda_{n,v}}\,\|\theta_P(c)\|_2
+\sqrt{J_{\mathrm c}}\,\Lambda_{n,v}
\right\}
\]
is measurable and has probability at least \(0.9925\) under \(P^{\otimes n}\). Moreover, if \(\Lambda_{n,v}=0\), then \(P^{\otimes n}\)-almost surely, simultaneously for every \(c\in[-1/2,1/2]\),
\[
W(c)=\|\theta_P(c)\|_2^2.
\]
\end{lemmav}

\begin{theoremv}{thm:original-record-attainable-rate}[Original-record attainable separation rate]
For every admissible public tuple \(v=(p,\alpha,\beta,\gamma)\) of \cref{obj:def:model}, use the explicit rule \(\phi^*_{n,v}\) and its tuning quantities from \cref{obj:def:explicit-score-ledger,obj:def:sharp-frontier-scales}. Write
\[
q=\frac{p-1}{p},\qquad S=\alpha+\beta,\qquad
t=\frac{2-p}{p-1},\qquad
E_0=\frac{2\gamma q}{2\gamma+q},\qquad
E_4=\frac{2S}{1+2\alpha+\beta/q+S/(2\gamma)},\qquad
E=\min\{E_0,E_4\}.
\]
Define the minimum smoothness \(s_0\), increment summability exponent \(\eta_0\), geometric-series factor \(A(x)\), and attainment constants by
\[
\begin{aligned}
s_0&=\min\{\alpha,\beta,\gamma\},&
\lambda_0&=\frac{(p-1)^2}{4p},&
\eta_0&=\frac{(p-1)(p+6)}{4p},\\
A(x)&=(1-2^{-x})^{-1}\quad(x>0),\\
B_*&=420+800\{A(\alpha)+A(\lambda_0)\},\\
D_*&=1+\{\exp(1)S\log 2\}^{-1},\\
L_*&=16+880A(s_0)+160D_*+960A(\alpha),\\
G_*&=\sqrt{64}\{2^{t/2}A(1/2)+A(\eta_0/2)\},\\
A_*&=\sqrt{2}\{24576+64L_*^2+16384+256G_*^2\},\\
C^{\mathrm{att}}_v&=16\{16+5B_*+1024\sqrt{A_*}\}.
\end{aligned}
\]
Define the fixed witness distance \(d_0\) and public attainment threshold \(N^{\mathrm{att}}_v\) by
\[
d_0=\frac{1}{8\sqrt{2}},
\qquad
N^{\mathrm{att}}_v
=\left\lceil\max\left\{4,
\left(\frac{2C^{\mathrm{att}}_v}{d_0}\right)^{1/E}\right\}\right\rceil.
\]
Let \(\lambda\) be Lebesgue probability measure on \([0,1]\). Define the heterogeneity distance \(d(P)\) and its model supremum \(D_v\) by
\[
d(P)=
\left\{\int_{[0,1]}
\left(\tau_P(x)-\int_{[0,1]}\tau_P(z)\,d\lambda(z)\right)^2
\,d\lambda(x)\right\}^{1/2},
\qquad
D_v=\sup_{P\in\mathcal M_v}d(P).
\]
Let \(\mathcal D_n\) be a sample of \(n\) independent original records with law \(P\), and let \(U\) be an independent uniform randomization variable on \([0,1]\). For a test \(\phi\), define its rejection expectation by
\[
\mathsf R_n(P,\phi)=\mathbb E[\phi(\mathcal D_n,U)].
\]

Then \(C^{\mathrm{att}}_v>0\), and the following guarantees hold:
\begin{itemize}
\item \textbf{(Size and separation.)}
For every integer \(n\ge2\),
\[
\mathsf R_n(P,\phi^*_{n,v})\le0.0075
\qquad\text{for every }P\in H_0(v),
\]
where the null is given in \cref{obj:def:null}. Whenever
\(C^{\mathrm{att}}_v n^{-E}<D_v\),
\[
1-\mathsf R_n(P,\phi^*_{n,v})\le0.0075
\qquad
\text{for every }P\in\mathcal M_v
\text{ with }d(P)\ge C^{\mathrm{att}}_v n^{-E}.
\]
The capped critical radius of \cref{obj:def:critical-radius} satisfies
\[
r_n^*(v)\le C^{\mathrm{att}}_v n^{-E}.
\]

\item \textbf{(Witness threshold.)}
For every integer \(n\ge N^{\mathrm{att}}_v\),
\[
C^{\mathrm{att}}_v n^{-E}<d_0.
\]

\item \textbf{(Tuning bounds.)}
For every integer \(n\ge4\), the coarse rank \(M\), fine rank \(K\), block size \(s\), main clipping cutoff \(T_0\), and increment cutoffs \(T_j\) from \cref{obj:def:explicit-score-ledger,obj:def:blocks} satisfy
\[
M\le2s,\qquad K\le2s^2,\qquad
1\le T_0\le2s,\qquad
1\le T_j\le T_0\quad(1\le j\le L).
\]

\item \textbf{(Compact uniformity.)}
For every compact set \(K\) of public tuples in the product topology whose members are all admissible under \cref{obj:def:model}, there exist real constants \(C,N,e\), with \(e>0\), such that, for every \(v\in K\),
\[
C^{\mathrm{att}}_v\le C,\qquad
N^{\mathrm{att}}_v\le N,\qquad e\le E(v).
\]
Here \(E(v)\) denotes the displayed minimum exponent evaluated at \(v\).
\end{itemize}

For every admissible smoothness tuple \(w=(\alpha,\beta,\gamma)\) of \cref{obj:def:bounded-model}, define the bounded separation exponent \(E_{\mathrm b}\), matched attainment multiplier \(C^{\mathrm{att}}_{(2,w)}\), and bounded-model distance supremum \(D_w^{\mathrm b}\) by
\[
E_{\mathrm b}
=\min\left\{\frac{2\gamma}{4\gamma+1},
\frac{2(\alpha+\beta)}
{1+2(\alpha+\beta)+(\alpha+\beta)/(2\gamma)}\right\},
\qquad
C^{\mathrm{att}}_{(2,w)}
=C^{\mathrm{att}}_{(2,\alpha,\beta,\gamma)},
\qquad
D_w^{\mathrm b}=\sup_{P\in\mathcal M^{\mathrm b}_w}d(P).
\]
The original-model exponent evaluated at \((2,\alpha,\beta,\gamma)\) equals \(E_{\mathrm b}\). The bounded-model rule \(\phi^{*,\mathrm b}_{n,w}\) from \cref{obj:def:sharp-frontier-scales}, obtained by evaluating the same ledger rule at \((2,\alpha,\beta,\gamma)\), satisfies, for every integer \(n\ge2\),
\[
\mathsf R_n(P,\phi^{*,\mathrm b}_{n,w})\le0.0075
\qquad\text{for every }P\in H_0^{\mathrm b}(w),
\]
with the bounded null as in \cref{obj:def:bounded-null}. Whenever
\(C^{\mathrm{att}}_{(2,w)}n^{-E_{\mathrm b}}<D_w^{\mathrm b}\),
\[
1-\mathsf R_n(P,\phi^{*,\mathrm b}_{n,w})\le0.0075
\qquad
\text{for every }P\in\mathcal M^{\mathrm b}_w
\text{ with }d(P)\ge C^{\mathrm{att}}_{(2,w)}n^{-E_{\mathrm b}}.
\]
The bounded capped critical radius of \cref{obj:def:bounded-radius} satisfies
\[
r_n^{*,\mathrm b}(w)
\le C^{\mathrm{att}}_{(2,w)}n^{-E_{\mathrm b}}.
\]
\end{theoremv}

\begin{lemmav}{lem:paired-tent-full-record-lower}[Paired-tent full-record lower bounds]
Let \(\lambda\) denote Lebesgue probability measure on \([0,1]\), and let \(Q_{a,P}(dy\mid x)\) denote the conditional outcome probability kernel under law \(P\) in treatment arm \(a\in\{0,1\}\). Define the heterogeneity distance, its model suprema, and the fixed witness distance by
\[
d(P)=
\left\{\int_0^1
\left(\tau_P(x)-\int_0^1\tau_P(z)\,\lambda(dz)\right)^2
\,\lambda(dx)\right\}^{1/2},
\qquad
D_v=\sup_{P\in\mathcal M_v}d(P),
\qquad
D_w^{\mathrm b}=\sup_{P\in\mathcal M^{\mathrm b}_w}d(P),
\qquad
d_0=\frac{1}{8\sqrt2},
\]
where the models are defined in \cref{obj:def:model,obj:def:bounded-model}.

Write a normalized finite probability prior as
\[
\pi=\sum_{j=1}^{k}\omega_j\delta_{P_j},
\qquad
\omega_j\ge0,
\qquad
\sum_{j=1}^{k}\omega_j=1.
\]
Define its positive-weight support and full original-record mixture by
\[
\operatorname{supp}_+(\pi)=\{P_j:\omega_j>0\},
\qquad
\mathbb P_{\pi,n}=\sum_{j=1}^{k}\omega_jP_j^{\otimes n}.
\]
For probability laws on the same measurable space, define total variation by
\[
\operatorname{TV}(P,Q)=
\sup_{E\ \mathrm{measurable}}|P(E)-Q(E)|.
\]

For every admissible moment and smoothness tuple \(v=(p,\alpha,\beta,\gamma)\) as in \cref{obj:def:model}, put
\[
q=\frac{p-1}{p},
\qquad
E_0=\frac{2\gamma q}{2\gamma+q},
\qquad
c^{\mathrm e}_v=\frac{4^{-\gamma}}{16\sqrt3}.
\]
There exists a sequence of mark magnitudes \(L:\mathbb N\to\mathbb R\) with
\[
\lim_{n\to\infty}L(n)=+\infty
\]
such that, for every integer \(n\ge2\), there exist a null law \(P_0\) and a normalized finite probability prior \(\pi_1\), with at least one strictly positive weight, satisfying:
\begin{itemize}
\item \textbf{(Null and propensity.)}
The law \(P_0\) belongs to \(H_0(v)\) from \cref{obj:def:null}, and
\[
e_{P_0}(x)=\frac12
\qquad\text{for every }x\in[0,1].
\]

\item \textbf{(Arm support and moments.)}
The magnitude \(L(n)\) is positive. For \(P=P_0\) and every \(P\in\operatorname{supp}_+(\pi_1)\), each arm \(a\in\{0,1\}\) satisfies, for \(\lambda\)-almost every \(x\),
\[
Q_{a,P}\bigl(\{0,-L(n),L(n)\}\mid x\bigr)=1,
\qquad
\int_{\mathbb R}|y|^p\,Q_{a,P}(dy\mid x)\le1.
\]

\item \textbf{(Alternative placement.)}
Every \(P\in\operatorname{supp}_+(\pi_1)\) belongs to \(\mathcal M_v\) from \cref{obj:def:model}, satisfies \(e_P(x)=1/2\) for every \(x\in[0,1]\), and obeys
\[
c^{\mathrm e}_v n^{-E_0}\le d(P)<d_0.
\]
Moreover, \(d_0\le D_v\).

\item \textbf{(Full-record comparison.)}
Define the alternative mixture by
\[
\overline P_1^{(n)}=\mathbb P_{\pi_1,n}.
\]
The same witnesses satisfy
\[
\operatorname{TV}\bigl(P_0^{\otimes n},\overline P_1^{(n)}\bigr)<\frac12.
\]
\end{itemize}
For every such \(n\), the supplied-propensity and original-record minimax testing risks in \cref{obj:def:oracle-risk,obj:def:testing-risk} satisfy
\[
B_n^{\mathrm e}\bigl(v,c^{\mathrm e}_v n^{-E_0}\bigr)\ge\frac25,
\qquad
B_n\bigl(v,c^{\mathrm e}_v n^{-E_0}\bigr)\ge\frac25.
\]

For every admissible smoothness tuple \(w=(\alpha,\beta,\gamma)\) as in \cref{obj:def:bounded-model}, put
\[
c_w^{\mathrm b}=\frac{4^{-\gamma}}{16\sqrt3}.
\]
For every integer \(n\ge2\), there exist a signed-binary null law \(P_0^{\mathrm b}\) and a normalized finite probability prior \(\pi_1^{\mathrm b}\), with at least one strictly positive weight, satisfying:
\begin{itemize}
\item \textbf{(Signed-binary null.)}
The law \(P_0^{\mathrm b}\) belongs to \(H_0^{\mathrm{bin}}(w)\) from \cref{obj:def:binary-null}, and
\[
e_{P_0^{\mathrm b}}(x)=\frac12
\qquad\text{for every }x\in[0,1].
\]

\item \textbf{(Signed-binary alternatives.)}
Every \(P\in\operatorname{supp}_+(\pi_1^{\mathrm b})\) belongs to \(\mathcal M^{\mathrm{bin}}_w\) from \cref{obj:def:binary-model}, satisfies \(e_P(x)=1/2\) for every \(x\in[0,1]\), and obeys
\[
c_w^{\mathrm b}n^{-2\gamma/(4\gamma+1)}
\le d(P)<d_0.
\]
Moreover, \(d_0\le D_w^{\mathrm b}\).

\item \textbf{(Full-record comparison.)}
Define the signed-binary alternative mixture by
\[
\overline P_{1,\mathrm b}^{(n)}
=\mathbb P_{\pi_1^{\mathrm b},n}.
\]
The same witnesses satisfy
\[
\operatorname{TV}\bigl((P_0^{\mathrm b})^{\otimes n},
\overline P_{1,\mathrm b}^{(n)}\bigr)<\frac12.
\]
\end{itemize}
For every such \(n\), the bounded-model minimax testing risk in \cref{obj:def:bounded-risk} satisfies
\[
B_n^{\mathrm b}\bigl(w,c_w^{\mathrm b}n^{-2\gamma/(4\gamma+1)}\bigr)
\ge\frac25.
\]
\end{lemmav}

\begin{definitionv}{def:marked-table}[Six-category record table \(\mathsf T\)]
The conditional original-record table is
\[
\begin{aligned}
\mathsf T\bigl(A=(1+\ell)/2,Y=hL\mid X=x\bigr)
&=\frac{\varepsilon}{4}
\{1+\ell\xi(x)+h\upsilon(x)+\ell h\zeta(x)\},
&&(\ell,h)\in\{-1,1\}^2,\\
\mathsf T\bigl(A=(1+\ell)/2,Y=0\mid X=x\bigr)
&=\frac{1-\varepsilon}{2}\{1+\ell\xi(x)\},
&&\ell\in\{-1,1\}.
\end{aligned}
\]
Here \(\varepsilon\) is the rare-mark probability, \(L>0\) is the mark magnitude, and \(\xi,\upsilon,\zeta\) may be Borel functions of the covariate. The table is evaluated on inputs for which every displayed probability is nonnegative. Likelihood products include all six record categories.
\end{definitionv}

\begin{definitionv}{def:frame-handle}[Finite-overlap frame \(f_i\)]
The finite-overlap frame at rank \(K=J_{\mathrm f}\), the fine histogram rank, consists of the functions
\[
f_i(x)=
\begin{cases}
\cos\!\bigl(\pi(Kx-i)/2\bigr),
&x\in[i/K,(i+1)/K],\\
\sin\!\bigl(\pi(Kx-(i-1))/2\bigr),
&i>0,\ x\in[(i-1)/K,i/K],\\
0,&\text{otherwise},
\end{cases}
\qquad x\in[0,1],\quad 0\le i\le K.
\]
For a coefficient vector
\[
a=(a_0,\ldots,a_K),
\]
the associated frame field is
\[
x\longmapsto\sum_{i=0}^{K}a_i f_i(x).
\]
\end{definitionv}

\begin{algorithmv}{def:copula-frame-priors}[Copula-frame priors \(\pi^{K,M}_{\nu;v,a,u,\varepsilon,L}\)]
Inputs are the public tuple \(v\), dyadic ranks \(K=J_{\mathrm f}\) and \(M=J_{\mathrm c}\) satisfying \(M\ge2\) and \(K\ge16M\), amplitudes \(0<a,u\le1/16\), a rare-mark probability \(0<\varepsilon<1\), and a mark magnitude \(L>0\).
\begin{enumerate}
\item Set the coupling amplitude and triangular tent to
\[
\kappa=\frac1{16},
\qquad
b_\triangle(t)=2\min\{t,1-t\},
\quad t\in[0,1].
\]
Draw independent uniform signs
\[
\Pr(\sigma_j=1)=\Pr(\sigma_j=-1)=\frac12,
\qquad 1\le j\le M/2.
\]
Define the opposite paired coarse tents by
\[
g_\sigma(x)=
\begin{cases}
\sigma_j b_\triangle(Mx-(2j-2)),
&x\in[(2j-2)/M,(2j-1)/M],\\
-\sigma_j b_\triangle(Mx-(2j-1)),
&x\in[(2j-1)/M,2j/M],
\end{cases}
\quad 1\le j\le M/2.
\]
Their values at coarse-cell boundaries are zero.
\item On each fine cell define the two nonzero frame coordinates by
\[
f_i(x)=\cos\!\bigl(\pi(Kx-i)/2\bigr),
\qquad
f_{i+1}(x)=\sin\!\bigl(\pi(Kx-i)/2\bigr),
\]
\[
x\in[i/K,(i+1)/K],
\qquad 0\le i<K.
\]
Set the other coordinates to zero and form the squared-frame interpolation
\[
\widetilde g_\sigma(x)
=
\sum_{i=0}^{K}g_\sigma(i/K)f_i(x)^2.
\]
\item For each branch \(\nu\in\{0,1\}\), draw coefficient pairs independently conditional on \(\sigma\), with
\[
\Pr_\nu(\lambda_i=l,\eta_i=h\mid\sigma)
=
\frac{1+\nu\kappa g_\sigma(i/K)lh}{4},
\qquad
l,h\in\{-1,1\},
\quad 0\le i\le K.
\]
\item For each coefficient draw, define the conditional-table coordinates
\[
\xi(x)=a\sum_{i=0}^{K}f_i(x)\lambda_i,
\qquad
\upsilon(x)=u\sum_{i=0}^{K}f_i(x)\eta_i,
\]
\[
t_\nu(x)=
-\frac{\nu au\kappa}{1-a^2}\widetilde g_\sigma(x),
\qquad
\zeta(x)=
\xi(x)\upsilon(x)+t_\nu(x)\bigl(1-\xi(x)^2\bigr).
\]
Let \(P\) be the original-record law with uniform covariate \(X\) and the normalized six-category marked conditional table evaluated at these coordinates, rare-mark probability \(\varepsilon\), and mark magnitude \(L\).
\item Output the finite prior obtained by pushing the coefficient distribution forward to original-record laws:
\[
\pi^{K,M}_{\nu;v,a,u,\varepsilon,L}
=
\mathcal L_\nu(P),
\qquad \nu\in\{0,1\}.
\]
\end{enumerate}
\end{algorithmv}

\begin{lemmav}{lem:copula-frame-legality}[Copula-frame legality and moments]
Consider an admissible tuple \(v=(p,\alpha,\beta,\gamma)\) as stipulated in \cref{obj:def:model}. Put
\[
q=\frac{p-1}{p},\qquad S=\alpha+\beta,\qquad
F(p)=\frac{2\alpha}{p-1}+\frac{\beta}{q}+\frac{S}{2\gamma}.
\]
Suppose the following conditions hold:
\begin{itemize}
\item \textbf{(Phase.)} \(F(p)<1\).
\item \textbf{(Dyadic ranks.)} \(K\) and \(M\) are powers of two, with \(16M\le K\) and \(2\le M\).
\item \textbf{(Rank compatibility.)} \(M\le K^{S/\gamma}\).
\end{itemize}
Define the construction amplitudes by
\[
a=\frac{K^{-\alpha}}{16},\qquad b=\frac{K^{-\beta}}{256}.
\]
Let \(\lambda\) denote Lebesgue probability measure on \([0,1]\), and write \(Q_{a,P}(dy\mid x)\) for the conditional outcome kernel in treatment arm \(a\) under \(P\). Define the population effect distance by
\[
d(P)=
\left(
\int_0^1
\left|\tau_P(x)-\int_0^1\tau_P(z)\,\lambda(dz)\right|^2
\,\lambda(dx)
\right)^{1/2},
\]
and let \(d_0\) be the fixed positive distance supplied by the legal model witness.

Use the copula-frame priors of \cref{obj:def:copula-frame-priors} with
\[
u=\frac1{16},\qquad
\varepsilon=(16b)^{1/q},\qquad L=\varepsilon^{-1/p},
\qquad
\pi_0=\pi^{K,M}_{0;v,a,u,\varepsilon,L},\qquad
\pi_1=\pi^{K,M}_{1;v,a,u,\varepsilon,L}.
\]
For every realization of the \(M/2\) coarse signs and the \(K+1\) fine coefficient-sign pairs, the null-branch law belongs to \(H_0(v)\) of \cref{obj:def:null}, and the alternative-branch law belongs to \(\mathcal M_v\) of \cref{obj:def:model}. Each alternative-branch law satisfies
\[
\tau_P(x)=-\frac{2ab\kappa}{1-a^2}\widetilde g_\sigma(x)
\quad\text{for every }x\in[0,1],
\qquad
\int_0^1\tau_P(x)\,\lambda(dx)=0,
\]
\[
2^{-17}K^{-S}\le d(P)\le 2^{-14}K^{-S},
\qquad d(P)<d_0.
\]
Every law in either branch satisfies
\[
\int_{\mathbb R}|y|^p\,Q_{a,P}(dy\mid x)=1
\quad\text{for every }a\in\{0,1\}\text{ and }x\in[0,1].
\]

For the bounded construction, let \(w=(\alpha,\beta,\gamma)\) be an admissible tuple as stipulated in \cref{obj:def:bounded-model}, set \(v=(2,\alpha,\beta,\gamma)\), and impose the same phase, dyadic-rank, and rank-compatibility conditions with \(p=2\). Retain the displayed amplitudes \(a,b\), and use the priors of \cref{obj:def:copula-frame-priors} with
\[
u=2b,\qquad \varepsilon=\frac12,\qquad L=1,
\qquad
\pi_0=\pi^{K,M}_{0;v,a,u,\varepsilon,L},\qquad
\pi_1=\pi^{K,M}_{1;v,a,u,\varepsilon,L}.
\]
For every realization of the coarse signs and fine coefficient-sign pairs, the null-branch law belongs to \(H_0^{\mathrm b}(w)\) of \cref{obj:def:bounded-null}, and the alternative-branch law belongs to \(\mathcal M^{\mathrm b}_w\) of \cref{obj:def:bounded-model}. Each alternative-branch law satisfies the displayed effect identity, centering identity, and distance bounds. Every law in either bounded branch satisfies
\[
\int_{\mathbb R}|y|^2\,Q_{a,P}(dy\mid x)=\frac12
\quad\text{for every }a\in\{0,1\}\text{ and }x\in[0,1].
\]
\end{lemmav}

\begin{lemmav}{lem:full-record-copula-component-bound}[Full-record copula component bound]
Let \(v=(p,\alpha,\beta,\gamma)\) be any real tuple, and suppose:
\begin{itemize}
\item \textbf{(Sample size and ranks.)} \(n,K,M\) are natural numbers, \(n\ge2\), \(K\) and \(M\) are powers of two, \(M\ge2\), \(16M\le K\), and \(2^{40}n\le K\).
\item \textbf{(Amplitudes.)} \(0<a\le1/16\) and \(0<u\le1/16\).
\item \textbf{(Mark parameters.)} \(0<\varepsilon<1\) and \(L>0\).
\end{itemize}
Write \(\pi_0\) and \(\pi_1\) for the two copula-frame priors
\[
\pi_\nu=\pi^{K,M}_{\nu;v,a,u,\varepsilon,L},
\qquad \nu\in\{0,1\},
\]
and define their complete-record mixtures by
\[
\mathbb P_\nu=\int P^{\otimes n}\,\pi_\nu(dP),
\qquad \nu\in\{0,1\}.
\]
Let \(\lambda\) denote Lebesgue probability measure on \([0,1]\). Total variation, denoted by \(\operatorname{TV}\), is the supremum of the absolute difference of event probabilities. The directed chi-square divergence \(\chi^2\) is the integral of the squared Radon--Nikodym density minus one, for probability laws with the stated absolute continuity.

There exist a common augmentation law \(\mu\), conditional label kernels \(P_0,P_1,Q\), and nonnegative measurable, integrable conditional budgets \(\mathcal A,\mathcal B\) with the following properties. The augmentation coordinate is
\[
\left((X_i)_{i=1}^n,(B_i)_{i=1}^n,\mathcal D_{\partial}\right),
\qquad B_i=\mathbf1\{Y_i\ne0\},
\]
where \(\mathcal D_{\partial}\) records the sampled coefficient pairs
\((\lambda_j,\eta_j)\) exactly at nodes satisfying
\[
jM\bmod K=0,\qquad 0\le j\le K.
\]
Its covariate marginal is \(\lambda^{\otimes n}\), and its flag marginal is
\(\operatorname{Bernoulli}(\varepsilon)^{\otimes n}\).

Each label consists of a binary treatment and a binary outcome sign, retaining the sign also when the mark flag is zero. Under branch \(\nu\), the joint law of the augmentation and these labels induced by the latent copula construction is \(\mu\otimes P_\nu\); the disclosure equals the boundary restriction of the sampled coefficients almost surely. Under each branch, the \(M/2\) coarse signs are independent fair signs, and their joint law with the augmentation is the product of their fair-sign law and \(\mu\).

For each branch, recovering the records by retaining \(X_i\) and the treatment label and setting the outcome to the signed magnitude \(L\) when \(B_i=1\), and to zero when \(B_i=0\), sends \(\mu\otimes P_\nu\) to \(\mathbb P_\nu\). Both augmented branch laws have augmentation marginal \(\mu\).

The intermediate augmented law \(\mu\otimes Q\) is a probability law. For \(\mu\)-almost every augmentation, \(P_0,P_1,Q\) are probability laws on the labels and satisfy
\[
Q\ll P_0,\qquad P_1\ll Q,\qquad
\mathcal A\ge0,\qquad \mathcal B\ge0,
\]
\[
1+\chi^2(Q\Vert P_0)\le \exp(\mathcal A),
\qquad
1+\chi^2(P_1\Vert Q)\le \exp(\mathcal B).
\]
Here all three conditional laws and both budgets are evaluated at the full augmentation triple. Their expectations satisfy
\[
\int \mathcal A\,d\mu
\le \frac{2^{38}a^4u^4\varepsilon^2 n^2}{K},
\qquad
\int \mathcal B\,d\mu
\le \frac{2^{56}a^4u^4\varepsilon^2 n^4}{MK^2}.
\]
If
\[
\frac{2^{38}a^4u^4\varepsilon^2 n^2}{K}
+
\frac{2^{56}a^4u^4\varepsilon^2 n^4}{MK^2}
\le 2^{-16},
\]
then
\[
\operatorname{TV}(\mathbb P_0,\mathbb P_1)<\frac18.
\]
\end{lemmav}

\begin{algorithmv}{def:converse-handle}[Branchwise converse priors \(\pi_{0,n,v},\pi_{1,n,v}\)]
Inputs are the sample size \(n\) and the public tuple \(v=(p,\alpha,\beta,\gamma)\); \(q=(p-1)/p\), \(F(p)\) is the exponent-comparison scalar, \(D_p\) is the rough-confounding exponent denominator, \(H_{\mathrm{fine}}\) is the fine-rank multiplier, and \(N_{\mathrm{low}}(v)\) is the rough-construction threshold.
\begin{enumerate}
\item If \(F(p)<1\) and \(n\ge N_{\mathrm{low}}(v)\), select
\[
K=2^{\left\lceil\log_2\!\left(H_{\mathrm{fine}}n^{2/D_p}\right)\right\rceil},
\qquad
M=2^{\left\lfloor\log_2\!\left(K^{(\alpha+\beta)/\gamma}/16\right)\right\rfloor},
\]
and set
\[
a=\frac{K^{-\alpha}}{16},
\qquad
u=\frac1{16},
\qquad
\varepsilon_{n,v}=16^{-1/q}K^{-\beta/q},
\qquad
L_{n,v}=\varepsilon_{n,v}^{-1/p}.
\]
Select the finite copula-frame priors
\[
\pi_{\nu,n,v}
=
\pi^{K,M}_{\nu;v,a,u,\varepsilon_{n,v},L_{n,v}},
\qquad \nu\in\{0,1\}.
\]
\item Otherwise, set
\[
N=2\left\lceil n^{2q/(2\gamma+q)}\right\rceil,
\qquad
h=\frac1N,
\qquad
\varepsilon_{n,v}=h^{\gamma/q},
\qquad
L_{n,v}=h^{-\gamma/(p-1)}.
\]
Select the opposite paired-tent null and alternative priors at this rank, cell width, rare-mark probability, and mark magnitude, denoted respectively by
\[
\pi_{0,n,v},\qquad \pi_{1,n,v}.
\]
\item Output the selected priors and their complete independent-record mixtures
\[
\bigl(\pi_{0,n,v},\pi_{1,n,v}\bigr),
\qquad
\mathbb P_{\nu,n,v}
=
\int P^{\otimes n}\,\pi_{\nu,n,v}(dP),
\quad \nu\in\{0,1\}.
\]
The selected construction satisfies the stipulated balancing, support-law membership, true null-mixture denominator, complete occupancy and transition accounting, and bounded-scale comparison properties.
\end{enumerate}
\end{algorithmv}

\begin{definitionv}{synth_1}[Bounded converse priors]
Fix \(v=(2,\alpha,\beta,\gamma)\in\mathcal V\), write \(w=(\alpha,\beta,\gamma)\), and let \(n\ge2\). Use the phase functional and public constants of \cref{obj:def:sharp-frontier-scales}. For \(\nu\in\{0,1\}\), define the bounded prior pair \(\pi^{\mathrm b}_{\nu,n,v}\) as follows.

If \(F(2)<1\) and \(n\ge N_{\mathrm{low}}(v)\), put \(S=\alpha+\beta\), let \(K\) be the least power of two at least \(H_{\mathrm{fine}}n^{2/D_2}\), and set
\[
 M=2^{\lfloor\log_2(K^{S/\gamma}/16)\rfloor},
 \qquad a=\frac{K^{-\alpha}}{16},
 \qquad b=\frac{K^{-\beta}}{256}.
\]
Define
\[
 \pi^{\mathrm b}_{\nu,n,v}
   =\pi^{K,M}_{\nu;v,a,\,2b,\,1/2,\,1},
 \qquad \nu\in\{0,1\},
\]
using the finite coefficient draws and record-law pushforward of \cref{obj:def:copula-frame-priors}. This is the bounded tuning of \cref{obj:lem:copula-frame-legality}: it uses the same ranks and amplitudes \(a,b\) as the rough branch of \cref{obj:def:converse-handle}, with outcome-coordinate amplitude \(u=2b\), mark probability \(\varepsilon=1/2\), and mark magnitude \(L=1\).

If \(F(2)\ge1\), or if \(F(2)<1\) and \(n<N_{\mathrm{low}}(v)\), take the signed-binary paired-tent package of \cref{obj:lem:paired-tent-full-record-lower} at sample size \(n\) and tuple \(w\). Writing \(P_0^{\mathrm b}\) for its selected null law and \(\pi_1^{\mathrm b}\) for its selected finite alternative prior, define
\[
 \pi^{\mathrm b}_{0,n,v}=\delta_{P_0^{\mathrm b}},
 \qquad
 \pi^{\mathrm b}_{1,n,v}=\pi_1^{\mathrm b}.
\]
This supplies the signed-binary counterpart of the paired-tent branch and small-sample fallback in \cref{obj:def:converse-handle}.

In either branch, the associated complete-record mixtures are
\[
 \mathbb P^{\mathrm b}_{\nu,n,v}
   =\int P^{\otimes n}\,\pi^{\mathrm b}_{\nu,n,v}(dP),
 \qquad \nu\in\{0,1\}.
\]
\end{definitionv}

\begin{lemmav}{lem:rough-full-record-sharp-lower}[Sharp rough-phase record lower bound]
Use the selected converse priors and complete-record mixtures of \cref{obj:def:converse-handle}, and the phase functional, exponents, multipliers, and scales of \cref{obj:def:sharp-frontier-scales}. Let \(\lambda\) denote Lebesgue probability measure on \([0,1]\). Define the heterogeneity distance, model supremum, and witness distance by
\[
d(P)=\left(\int_{[0,1]}
\left(\tau_P(x)-\int_{[0,1]}\tau_P(z)\,d\lambda(z)\right)^2
\,d\lambda(x)\right)^{1/2},
\qquad
D_v=\sup_{P\in\mathcal M_v}d(P),
\qquad
d_0=\frac{1}{8\sqrt2}.
\]
For probability laws, \(\operatorname{TV}\) denotes the supremum of the absolute differences between their probabilities of the same measurable event.

Suppose that:
\begin{itemize}
\item \textbf{(Admissibility.)} \(v=(p,\alpha,\beta,\gamma)\) is admissible as in \cref{obj:def:model}.
\item \textbf{(Rough phase.)} \(F(p)<1\).
\item \textbf{(Sample size.)} \(n\) is an integer with \(n\ge2\).
\end{itemize}
The finite priors \(\pi_{0,n,v}\) and \(\pi_{1,n,v}\) have nonnegative weights summing to one. Every positive-weight null support law belongs to \(H_0(v)\) of \cref{obj:def:null}, and every positive-weight alternative support law satisfies
\[
P\in\mathcal M_v,
\qquad
c_vn^{-E_4(p)}\le d(P)<d_0.
\]
Moreover,
\[
d_0\le D_v,
\qquad
\operatorname{TV}(\mathbb P_{0,n,v},\mathbb P_{1,n,v})<\frac12,
\qquad
B_n\!\left(v,c_vn^{-E_4(p)}\right)\ge\frac25,
\]
with testing risk as in \cref{obj:def:testing-risk}.

If \(n\ge N_{\mathrm{low}}(v)\), the selected fine and coarse ranks \(K,M\) of \cref{obj:def:converse-handle} satisfy the rank conditions of \cref{obj:lem:copula-frame-legality}. In particular, they are powers of two and satisfy
\[
16M\le K,
\qquad
2\le M\le K^{S/\gamma},
\qquad
2^{40}n\le K.
\]
With the selected rare-mark probability \(\varepsilon_{n,v}\) of \cref{obj:def:converse-handle}, the deterministic budget bounds are
\[
\begin{aligned}
\frac{2^{38}(K^{-\alpha}/16)^4(1/16)^4
      \varepsilon_{n,v}^{\,2}n^2}{K}
&\le \frac{2^{-10}}{H_{\mathrm{fine}}},\\
\frac{2^{56}(K^{-\alpha}/16)^4(1/16)^4
      \varepsilon_{n,v}^{\,2}n^4}{MK^2}
&\le \frac{2^{13}}{H_{\mathrm{fine}}^2}.
\end{aligned}
\]
The same complete-record mixtures then satisfy
\[
\operatorname{TV}(\mathbb P_{0,n,v},\mathbb P_{1,n,v})<\frac18.
\]

For the bounded experiment, suppose that:
\begin{itemize}
\item \textbf{(Admissibility.)} \(w=(\alpha,\beta,\gamma)\) is admissible as in \cref{obj:def:bounded-model}; write \(v=(2,\alpha,\beta,\gamma)\).
\item \textbf{(Rough phase.)} \(F(2)<1\).
\item \textbf{(Sample size.)} \(n\) is an integer with \(n\ge2\).
\end{itemize}
Let \(\pi^{\mathrm b}_{0,n,v}\) and \(\pi^{\mathrm b}_{1,n,v}\) be the selected bounded converse priors of \cref{obj:synth_1}, and let
\(\mathbb P^{\mathrm b}_{0,n,v}\) and \(\mathbb P^{\mathrm b}_{1,n,v}\) denote their respective complete-record mixtures. Both finite priors have nonnegative weights summing to one. Every positive-weight null support law belongs to \(H_0^{\mathrm b}(w)\) of \cref{obj:def:bounded-null}, and every positive-weight alternative support law satisfies
\[
P\in\mathcal M^{\mathrm b}_w,
\qquad
c_w^{\mathrm b}n^{-E_4(2)}\le d(P)<d_0.
\]
Furthermore,
\[
\operatorname{TV}
\bigl(\mathbb P^{\mathrm b}_{0,n,v},\mathbb P^{\mathrm b}_{1,n,v}\bigr)
<\frac12,
\qquad
B_n^{\mathrm b}\!\left(w,c_w^{\mathrm b}n^{-E_4(2)}\right)\ge\frac25,
\]
with bounded testing risk as in \cref{obj:def:bounded-risk}.
\end{lemmav}

\begin{theoremv}{thm:tail-essential-full-record-converse}[Tail essentiality and bounded witnesses]
Fix an admissible tuple \(v=(p,\alpha,\beta,\gamma)\) as in \cref{obj:def:model}, and write \(w=(\alpha,\beta,\gamma)\). Use the scales, constants, thresholds, and total tests of \cref{obj:def:sharp-frontier-scales}, and the selected priors, mark magnitudes, and complete-record mixtures of \cref{obj:def:converse-handle}.

Let \(\lambda\) denote Lebesgue probability measure on \([0,1]\). Define the heterogeneity distance and its model supremum by
\[
d(P)=
\left\{\int_{[0,1]}
\left(\tau_P(x)-\int_{[0,1]}\tau_P(z)\,\lambda(dz)\right)^2
\,\lambda(dx)\right\}^{1/2},
\qquad
D_v=\sup_{P\in\mathcal M_v}d(P).
\]
For probability laws, \(\operatorname{TV}\) denotes the supremum of the absolute differences of their probabilities over measurable events.

A normalized finite prior on original-record laws is
\[
\pi=\sum_{j=1}^{k}\omega_j\delta_{P_j},
\qquad
k\ge1,\quad \omega_j\ge0,\quad \sum_{j=1}^{k}\omega_j=1.
\]
Its positive-weight support and complete-record mixture are defined by
\[
\operatorname{supp}_+(\pi)=\{P_j:\omega_j>0\},
\qquad
\mathbb P_{\pi,n}=\sum_{j=1}^{k}\omega_jP_j^{\otimes n}.
\]

For every integer \(n\ge2\), both \(\pi_{0,n,v}\) and \(\pi_{1,n,v}\) are normalized finite priors, and
\[
\begin{gathered}
\operatorname{supp}_+(\pi_{0,n,v})\subseteq H_0(v),\\
\operatorname{supp}_+(\pi_{1,n,v})
\subseteq
\{P\in\mathcal M_v:c_v\rho_n(v)\le d(P)<D_v\},\\
0<c_v\rho_n(v)<D_v,\\
\operatorname{TV}(\mathbb P_{0,n,v},\mathbb P_{1,n,v})<\frac12,
\qquad
B_n(v,c_v\rho_n(v))\ge\frac25,
\end{gathered}
\]
where the classes and testing risk are those of
\cref{obj:def:null,obj:def:model,obj:def:testing-risk}.
For every positive-weight law in either prior, each conditional outcome arm assigns probability one to
\(\{0,-L_{n,v},L_{n,v}\}\) for \(\lambda\)-almost every covariate value. Moreover,
\[
L_{n,v}\longrightarrow\infty
\qquad(n\longrightarrow\infty).
\]

Use the matched bounded model \(\mathcal M^{\mathrm b}_w\) and null
\(H_0^{\mathrm b}(w)\) of
\cref{obj:def:bounded-model,obj:def:bounded-null}.
If \(p<2\), then
\[
\frac{\rho_n(v)}{\rho_n^{\mathrm b}(w)}
\longrightarrow\infty
\qquad(n\longrightarrow\infty).
\]
For each such fixed \(v\), there exists an integer \(N(v)\ge2\) such that, for every integer \(n\ge N(v)\) and every pair of normalized finite priors \(\pi_0,\pi_1\),
\[
\left.
\begin{gathered}
\operatorname{supp}_+(\pi_0)\subseteq H_0^{\mathrm b}(w),\\
\operatorname{supp}_+(\pi_1)\subseteq
\{P\in\mathcal M^{\mathrm b}_w:c_v\rho_n(v)\le d(P)\}
\end{gathered}
\right\}
\quad\Longrightarrow\quad
\operatorname{TV}(\mathbb P_{\pi_0,n},\mathbb P_{\pi_1,n})
\ge\frac12.
\]

If \(p=2\), the following conclusions hold for every integer \(n\ge2\).
Let \(\pi^{\mathrm b}_{\nu,n,v}\), for \(\nu\in\{0,1\}\), denote the selected bounded witness priors: the bounded rough priors of
\cref{obj:lem:rough-full-record-sharp-lower} when
\(F(2)<1\) and \(n\ge N_{\mathrm{low}}(v)\), and the signed-binary paired-tent priors of
\cref{obj:lem:paired-tent-full-record-lower} otherwise. Define
\[
\mathbb P^{\mathrm b}_{\nu,n,v}
=\mathbb P_{\pi^{\mathrm b}_{\nu,n,v},n},
\qquad \nu\in\{0,1\}.
\]
Both priors are normalized and have nonempty positive-weight supports. Their positive-weight masses satisfy
\[
\sum_{P\in\operatorname{supp}_+(\pi^{\mathrm b}_{\nu,n,v})}
\pi^{\mathrm b}_{\nu,n,v}(\{P\})=1,
\qquad \nu\in\{0,1\}.
\]
Let \(d_0>0\) be the fixed witness distance of
\cref{obj:lem:causal-nonempty}. The support, total-variation, and risk conclusions are
\[
\begin{gathered}
\operatorname{supp}_+(\pi^{\mathrm b}_{0,n,v})
\subseteq H_0^{\mathrm b}(w),
\qquad
\operatorname{supp}_+(\pi^{\mathrm b}_{0,n,v})
\subseteq H_0(v),\\
\operatorname{supp}_+(\pi^{\mathrm b}_{1,n,v})
\subseteq
\{P\in\mathcal M^{\mathrm b}_w:c_v\rho_n(v)\le d(P)<d_0\},\\
\operatorname{supp}_+(\pi^{\mathrm b}_{1,n,v})
\subseteq
\{P\in\mathcal M_v:c_v\rho_n(v)\le d(P)\},\\
\operatorname{TV}(\mathbb P^{\mathrm b}_{0,n,v},
                  \mathbb P^{\mathrm b}_{1,n,v})<\frac12,\\
\frac25\le B_n^{\mathrm b}(w,c_v\rho_n(v))
\le B_n(v,c_v\rho_n(v)),
\end{gathered}
\]
with bounded risk defined in \cref{obj:def:bounded-risk}.

Define the bounded-model distance supremum by
\[
D_w^{\mathrm b}=\sup_{P\in\mathcal M^{\mathrm b}_w}d(P).
\]
On this face, the scales, multipliers, thresholds, and distance suprema agree:
\[
\begin{gathered}
\rho_n(v)=\rho_n^{\mathrm b}(w),
\qquad c_v=c_w^{\mathrm b},
\qquad C_v=C_w^{\mathrm b},\\
N_v=N_w^{\mathrm b},
\qquad D_v=D_w^{\mathrm b}.
\end{gathered}
\]
The capped critical radii of
\cref{obj:def:critical-radius,obj:def:bounded-radius,obj:def:binary-radius}
satisfy
\[
\begin{gathered}
c_v\rho_n(v)\le r_n^*(v)\le C_v\rho_n(v),\\
c_w^{\mathrm b}\rho_n^{\mathrm b}(w)
\le r_n^{*,\mathrm b}(w)=r_n^{*,\mathrm{bin}}(w)
\le C_w^{\mathrm b}\rho_n^{\mathrm b}(w).
\end{gathered}
\]

Let \(\mathsf R_n(P,\phi)\) denote the expected rejection probability under \(n\) independent original records and the public randomization of
\cref{obj:def:original-record-experiment}. The total tests satisfy, for every law in the indicated null class,
\[
\begin{aligned}
\mathsf R_n(P,\phi^*_{n,v})&\le0.0075
&&\quad(P\in H_0(v)),\\
\mathsf R_n(P,\phi^{*,\mathrm b}_{n,w})&\le0.0075
&&\quad(P\in H_0^{\mathrm b}(w)).
\end{aligned}
\]
Whenever \(C_v\rho_n(v)<D_v\), both separated alternative classes are nonempty:
\[
\begin{gathered}
\{P\in\mathcal M_v:C_v\rho_n(v)\le d(P)\}\ne\varnothing,\\
\{P\in\mathcal M^{\mathrm b}_w:
C_w^{\mathrm b}\rho_n^{\mathrm b}(w)\le d(P)\}\ne\varnothing.
\end{gathered}
\]
Under this same condition, every law in the corresponding separated class satisfies
\[
\begin{aligned}
1-\mathsf R_n(P,\phi^*_{n,v})&\le0.0075
&&\quad(P\in\mathcal M_v,\ C_v\rho_n(v)\le d(P)),\\
1-\mathsf R_n(P,\phi^{*,\mathrm b}_{n,w})&\le0.0075
&&\quad(P\in\mathcal M^{\mathrm b}_w,\\
C_w^{\mathrm b}\rho_n^{\mathrm b}(w)\le d(P)).
\end{aligned}
\]
For \(n\ge N_v\), the common upper separation satisfies
\[
\begin{gathered}
C_v\rho_n(v)\le\frac{d_0}{2},\\
C_v\rho_n(v)<D_v,
\qquad
C_w^{\mathrm b}\rho_n^{\mathrm b}(w)<D_w^{\mathrm b}.
\end{gathered}
\]
At \(n=2,3\), both total tests have zero rejection probability at every sample:
\[
\phi^*_{n,v}(z)=0,
\qquad
\phi^{*,\mathrm b}_{n,w}(z)=0
\qquad\text{for every }z.
\]
\end{theoremv}

\begin{theoremv}{thm:heavy-tail-comparison}[Heavy-tail critical-radius comparison]
Let the admissible exponent domains be
\[
\begin{aligned}
\mathcal V
&=\{(p,\alpha,\beta,\gamma)\in\mathbb R^4:
1<p\le2,\ 0<\alpha\le1,\ 0<\beta\le1,\ 1/4\le\gamma\le1\},\\
\mathcal W
&=\{(\alpha,\beta,\gamma)\in\mathbb R^3:
0<\alpha\le1,\ 0<\beta\le1,\ 1/4\le\gamma\le1\}.
\end{aligned}
\]
For \(v=(p,\alpha,\beta,\gamma)\in\mathcal V\), write
\(w=(\alpha,\beta,\gamma)\). Use the public scales, multipliers, thresholds, and original-record tests defined in \cref{obj:def:sharp-frontier-scales}, and the models, null classes, risks, and capped radii defined in \cref{obj:def:model,obj:def:null,obj:def:bounded-model,obj:def:bounded-null,obj:def:testing-risk,obj:def:bounded-risk,obj:def:critical-radius,obj:def:bounded-radius,obj:def:binary-radius}. The matched radius ratio from \cref{obj:def:comparison-handle} is
\[
\Delta_n(v)=\frac{r_n^*(v)}{r_n^{*,\mathrm b}(w)}.
\]
Here \(\lambda\) is Lebesgue probability measure on \([0,1]\), and the population distance of the mean effect from its average is
\[
d(P)=
\left(
\int_0^1
\left|\tau_P(x)-\int_0^1\tau_P(z)\,d\lambda(z)\right|^2
\,d\lambda(x)
\right)^{1/2}.
\]
The maximal distances over the respective models are
\[
D_v=\sup_{P\in\mathcal M_v}d(P),
\qquad
D_w^{\mathrm b}=\sup_{P\in\mathcal M^{\mathrm b}_w}d(P),
\]
and \(d_0>0\) is the fixed witness distance supplied by \cref{obj:lem:causal-nonempty}. For the original-record experiment in \cref{obj:def:original-record-experiment}, write the rejection expectation as
\[
\mathsf R_n(P,\phi)=\mathbb E_{\mathsf S_{n,P}}[\phi].
\]

The following assertions hold.
\begin{itemize}
\item \textbf{(Matched radii and lower separations.)}
For every \(v\in\mathcal V\), \(c_v>0\) and \(C_v>0\). For every integer \(n\ge2\), both scales are positive and
\[
\begin{aligned}
c_v\rho_n(v)
&\le r_n^*(v)\le C_v\rho_n(v),\\
c_w^{\mathrm b}\rho_n^{\mathrm b}(w)
&\le r_n^{*,\mathrm b}(w)
=r_n^{*,\mathrm{bin}}(w)
\le C_w^{\mathrm b}\rho_n^{\mathrm b}(w).
\end{aligned}
\]
The lower separations satisfy
\[
0<c_v\rho_n(v)<D_v,
\qquad
0<c_w^{\mathrm b}\rho_n^{\mathrm b}(w)<D_w^{\mathrm b},
\]
and the corresponding minimax type-II errors satisfy
\[
B_n\!\left(v,c_v\rho_n(v)\right)>\frac1{10},
\qquad
B_n^{\mathrm b}\!\left(w,c_w^{\mathrm b}\rho_n^{\mathrm b}(w)\right)>\frac1{10}.
\]

\item \textbf{(Size and conditional power.)}
For every \(v\in\mathcal V\) and every integer \(n\ge2\),
\[
\begin{aligned}
\mathsf R_n(P,\phi^*_{n,v})&\le\frac1{10}
&&\text{for every }P\in H_0(v),\\
\mathsf R_n(P,\phi^{*,\mathrm b}_{n,w})&\le\frac1{10}
&&\text{for every }P\in H_0^{\mathrm b}(w).
\end{aligned}
\]
If \(C_v\rho_n(v)<D_v\), then
\[
1-\mathsf R_n(P,\phi^*_{n,v})\le\frac1{10}
\]
for every \(P\in\mathcal M_v\) satisfying \(d(P)\ge C_v\rho_n(v)\).
If \(C_w^{\mathrm b}\rho_n^{\mathrm b}(w)<D_w^{\mathrm b}\), then
\[
1-\mathsf R_n(P,\phi^{*,\mathrm b}_{n,w})\le\frac1{10}
\]
for every \(P\in\mathcal M^{\mathrm b}_w\) satisfying
\(d(P)\ge C_w^{\mathrm b}\rho_n^{\mathrm b}(w)\).

\item \textbf{(Vanishing scales and public thresholds.)}
For every \(v\in\mathcal V\),
\[
\rho_n(v)\longrightarrow0,
\qquad
\rho_n^{\mathrm b}(w)\longrightarrow0
\quad\text{as }n\longrightarrow\infty.
\]
For every natural number \(n\),
\[
\begin{aligned}
n\ge N_v
&\ \Longrightarrow\ C_v\rho_n(v)<d_0,\\
n\ge N_w^{\mathrm b}
&\ \Longrightarrow\ C_w^{\mathrm b}\rho_n^{\mathrm b}(w)<d_0.
\end{aligned}
\]

\item \textbf{(Phase identities and radius comparison.)}
Use the supplied-propensity exponent \(E_0(p)\), rough-confounding exponent \(E_4(p)\), phase-comparison scalar \(F(p)\), and denominator \(D_p\) from \cref{obj:def:sharp-frontier-scales}, with the smoothness tuple \(w\) fixed. With
\[
q=\frac{p-1}{p},
\qquad
E(p)=\min\{E_0(p),E_4(p)\},
\]
every \(v\in\mathcal V\) satisfies
\[
E_4(p)-E_0(p)
=
\frac{2\gamma q}{(2\gamma+q)D_p}\bigl(F(p)-1\bigr),
\]
and
\[
F(p)\ge1\ \Longrightarrow\ E(p)=E_0(p),
\qquad
F(p)<1\ \Longrightarrow\ E(p)=E_4(p).
\]
With
\[
D_2=1+2S+\frac{S}{2\gamma},
\qquad
F_2=2S+\frac{S}{2\gamma},
\]
the exponent difference, with evaluation at \(2\) preserving \(w\), is
\[
E(2)-E(p)=
\begin{cases}
\displaystyle
\frac{4\gamma^2(2-p)}{(4\gamma+1)(2\gamma p+p-1)},
&F_2\ge1,\\[6pt]
\displaystyle E_4(2)-E_0(p),
&F_2<1\le F(p),\\[4pt]
\displaystyle
\frac{2S\beta(2-p)}{(p-1)D_pD_2},
&F(p)<1.
\end{cases}
\]
For every integer \(n\ge2\),
\[
\frac{c_v}{C_w^{\mathrm b}}\,n^{E(2)-E(p)}
\le\Delta_n(v)
\le
\frac{C_v}{c_w^{\mathrm b}}\,n^{E(2)-E(p)},
\qquad
1\le\Delta_n(v).
\]

\item \textbf{(Strict and equality regions.)}
The strict-enlargement and bounded-ratio regions of \cref{obj:def:tail-region,obj:def:equality-region} are exactly
\[
\mathcal T=\{v\in\mathcal V:p<2\},
\qquad
\mathcal E=\{v\in\mathcal V:p=2\}.
\]
There exists a set \(U\subseteq\mathbb R^4\), open in the product topology of the four public exponents, such that
\[
U\cap\mathcal V\ne\varnothing,
\qquad
U\cap\mathcal V\subseteq\mathcal T.
\]
For every \(v\in\mathcal V\) with \(p<2\),
\[
E(p)<E(2),
\qquad
\frac{\rho_n(v)}{\rho_n^{\mathrm b}(w)}
\longrightarrow\infty
\quad\text{as }n\longrightarrow\infty.
\]

\item \textbf{(Tail-essential complete-record priors.)}
For every \(v\in\mathcal V\) with \(p<2\), the selected mark magnitudes from \cref{obj:def:converse-handle} satisfy
\[
L_{n,v}\longrightarrow\infty
\quad\text{as }n\longrightarrow\infty.
\]
For every integer \(n\ge2\), the selected priors
\(\pi_{0,n,v}\) and \(\pi_{1,n,v}\) from \cref{obj:def:converse-handle} are normalized finite probability priors. Every law receiving positive null-prior weight belongs to \(H_0(v)\), and every law receiving positive alternative-prior weight belongs to \(\mathcal M_v\) and satisfies
\[
c_v\rho_n(v)\le d(P)<D_v.
\]
For every law receiving positive weight in either prior, each conditional arm law assigns probability one to the three-point mark set: writing \(Q_{a,P}(dy\mid x)\) for the conditional outcome probability kernel in arm \(a\),
\[
Q_{a,P}\!\left(\{0,-L_{n,v},L_{n,v}\}\mid x\right)=1
\quad
\text{for }a\in\{0,1\}\text{ and }\lambda\text{-almost every }x.
\]
The complete-record mixtures defined in \cref{obj:def:converse-handle} satisfy
\[
\operatorname{TV}\!\left(\mathbb P_{0,n,v},\mathbb P_{1,n,v}\right)<\frac12,
\]
where total variation is the supremum of the absolute difference in event probabilities.

For each such \(v\), there exists a natural-number threshold \(N(v)\) such that, for every \(n\ge N(v)\), every pair of normalized finite probability priors supported respectively on
\[
H_0^{\mathrm b}(w)
\quad\text{and}\quad
\{P\in\mathcal M^{\mathrm b}_w:c_v\rho_n(v)\le d(P)\}
\]
has complete-record mixture total variation at least \(1/2\), in the sense of \cref{obj:thm:tail-essential-full-record-converse}.

\item \textbf{(Compact uniformity.)}
For every compact set of admissible public tuples, there exist real constants \(c,C,N\), with \(c>0\), such that
\[
c\le c_v,\qquad C_v\le C,\qquad N_v\le N
\]
for every tuple in that set. Likewise, for every compact set of admissible smoothness tuples, there exist real constants \(c,C,N\), with \(c>0\), such that
\[
c\le c_w^{\mathrm b},
\qquad
C_w^{\mathrm b}\le C,
\qquad
N_w^{\mathrm b}\le N
\]
for every smoothness tuple in that set.
\end{itemize}
\end{theoremv}

\begin{theoremv}{thm:oracle-frontier-broad-class}[Supplied-propensity broader-class frontier]
For every admissible public tuple \(v=(p,\alpha,\beta,\gamma)\) of \cref{obj:def:model}, consider the broader model and its constant-effect null in \cref{obj:def:oracle-broad-model,obj:def:oracle-broad-null}, with testing risk and capped critical radius as in \cref{obj:def:oracle-broad-risk,obj:def:oracle-broad-radius}. Put
\[
q=\frac{p-1}{p},\qquad
E_0=\frac{2\gamma q}{2\gamma+q},\qquad
\rho_n^{\mathrm e}(v)=n^{-E_0},\qquad
d_0=\frac{1}{8\sqrt2},
\]
and define the lower multiplier, upper multiplier, and public sample-size threshold by
\[
c_v^{\mathrm e}=\frac{4^{-\gamma}}{16\sqrt3},\qquad
C_v^{\mathrm e}=3\bigl(120+128\sqrt{160\sqrt2}\bigr),\qquad
N_v^{\mathrm e}
=\left\lceil\max\left\{2,
\left(\frac{2C_v^{\mathrm e}}{d_0}\right)^{1/E_0}
\right\}\right\rceil .
\]
Let \(\lambda\) denote Lebesgue probability measure on \([0,1]\), and write the heterogeneity distance as
\[
d(P)=
\left\{\int_{[0,1]}
\left(\tau_P(x)-\int_{[0,1]}\tau_P(z)\,\lambda(dz)\right)^2
\,\lambda(dx)\right\}^{1/2}.
\]
Then
\[
d_0\le D_v^{\mathrm{e,br}}\le\frac12,\qquad
\frac{1}{64\sqrt3}\le c_v^{\mathrm e},\qquad
0<C_v^{\mathrm e}.
\]

For every integer \(n\ge2\), the following guarantees hold.
\begin{itemize}
\item \textbf{(Critical radius.)}
\[
c_v^{\mathrm e}\rho_n^{\mathrm e}(v)
\le r_n^{*,\mathrm{e,br}}(v)
\le C_v^{\mathrm e}\rho_n^{\mathrm e}(v),
\qquad
B_n^{\mathrm{e,br}}
\bigl(v,c_v^{\mathrm e}\rho_n^{\mathrm e}(v)\bigr)
\ge\frac25.
\]

\item \textbf{(Calibration and power.)}
The supplied-propensity test \(\phi^{*,\mathrm e}_{n,v}\) of
\cref{obj:def:oracle-handle} satisfies
\[
\mathsf R_n^{\mathrm e}
(P,\phi^{*,\mathrm e}_{n,v})\le0.01
\qquad
\text{for every }P\in H_0^{\mathrm{e,br}}(v).
\]
Here \(\mathsf R_n^{\mathrm e}(P,\phi^{\mathrm e})\) is the expected rejection probability when the test receives the entire true continuous propensity \(e_P\), together with the independent original-record sample and public randomization. If
\[
C_v^{\mathrm e}\rho_n^{\mathrm e}(v)<D_v^{\mathrm{e,br}},
\]
then every \(P\in\mathcal M_v^{\mathrm{e,br}}\) with
\(d(P)\ge C_v^{\mathrm e}\rho_n^{\mathrm e}(v)\) satisfies
\[
1-\mathsf R_n^{\mathrm e}
(P,\phi^{*,\mathrm e}_{n,v})\le0.01.
\]

\item \textbf{(Finite lower-bound witnesses.)}
There exist a continuous function \(f\), a null law \(P_0\), and a normalized finite probability prior
\[
f(x)=\frac12\quad(x\in[0,1]),\qquad
\pi_1=\sum_{j=1}^{k}\omega_j\delta_{P_j},\qquad
\omega_j\ge0,\qquad
\sum_{j=1}^{k}\omega_j=1,
\]
with at least one strictly positive weight, such that
\[
P_0\in H_0(v)\cap H_0^{\mathrm{e,br}}(v),
\qquad e_{P_0}=f.
\]
Here \(H_0(v)\) is the original-model null of \cref{obj:def:null}.
For every \(j\) with \(\omega_j>0\),
\[
P_j\in\mathcal M_v\cap\mathcal M_v^{\mathrm{e,br}},
\qquad e_{P_j}=f,\qquad
c_v^{\mathrm e}\rho_n^{\mathrm e}(v)\le d(P_j)<d_0,
\]
where \(\mathcal M_v\) is defined in \cref{obj:def:model}.
Define the complete-record alternative mixture by
\[
\overline P_1^{(n)}
=\sum_{j=1}^{k}\omega_j P_j^{\otimes n}.
\]
With total variation defined as the supremum of absolute event-probability differences,
\[
\operatorname{TV}
\bigl(P_0^{\otimes n},\overline P_1^{(n)}\bigr)<\frac12.
\]
For every allowed randomized supplied-propensity test \(\phi^{\mathrm e}\) satisfying
\[
\mathsf R_n^{\mathrm e}(P,\phi^{\mathrm e})\le\frac1{10}
\qquad
\text{for every }P\in H_0^{\mathrm{e,br}}(v),
\]
there is an index \(j\) with \(\omega_j>0\) such that
\[
1-\mathsf R_n^{\mathrm e}(P_j,\phi^{\mathrm e})
\ge\frac25.
\]
\end{itemize}

For every integer \(n\ge N_v^{\mathrm e}\),
\[
C_v^{\mathrm e}\rho_n^{\mathrm e}(v)\le\frac{d_0}{2}.
\]

For each \(P\in\mathcal M_v^{\mathrm{e,br}}\), let
\(Q_{a,P}(dy\mid x)\) denote the conditional outcome probability kernel in treatment arm \(a\). The untruncated inverse-propensity score has conditional mean equal to the canonical effect: for \(\lambda\)-almost every \(x\),
\[
e_P(x)\int_{\mathbb R}\frac{y}{e_P(x)}\,Q_{1,P}(dy\mid x)
+
\bigl(1-e_P(x)\bigr)
\int_{\mathbb R}\frac{-y}{1-e_P(x)}\,Q_{0,P}(dy\mid x)
=\tau_P(x).
\]

The constants and exponent are locally uniform on compact subsets of the admissible public parameter domain, equipped with the product topology on \((p,\alpha,\beta,\gamma)\).
\end{theoremv}

\begin{theoremv}{thm:oracle-frontier-resolved}[Supplied-propensity frontier and attainment]
Let \(\mathcal V\) denote the admissible domain of public tuples \(v=(p,\alpha,\beta,\gamma)\) in \cref{obj:def:model}, equipped with the topology inherited from \(\mathbb R^4\). For \(v\in\mathcal V\), define
\[
q=\frac{p-1}{p},\qquad
E_0=\frac{2\gamma q}{2\gamma+q},\qquad
\rho_n^{\mathrm e}(v)=n^{-E_0},\qquad
c^{\mathrm e}_v=\frac{4^{-\gamma}}{16\sqrt3},\qquad
C^{\mathrm e}_v=3\left(120+128\sqrt{160\sqrt2}\right).
\]
Use the supplied-propensity rule \(\phi^{*,\mathrm e}_{n,v}\) of \cref{obj:def:oracle-handle}, and the risk and capped radius of \cref{obj:def:oracle-risk,obj:def:oracle-radius}. Write \(\mathsf R_n^{\mathrm e}(P,\phi^{\mathrm e})\) for the expected rejection probability when the entire true continuous propensity \(e_P\) is supplied to the test.

Let \(\lambda\) be Lebesgue probability measure on \([0,1]\), and define
\[
d(P)=
\left\{\int_0^1
\left(\tau_P(x)-\int_0^1\tau_P(z)\,d\lambda(z)\right)^2
\,d\lambda(x)\right\}^{1/2},
\qquad
D_v=\sup_{P\in\mathcal M_v}d(P).
\]
Let \(d_0>0\) be the fixed distance supplied by the legal model witness, and set
\[
N^{\mathrm e}_v=
\left\lceil
\max\left\{2,\left(\frac{2C^{\mathrm e}_v}{d_0}\right)^{1/E_0}\right\}
\right\rceil.
\]
Then
\[
\frac{1}{64\sqrt3}\le c^{\mathrm e}_v,
\qquad
0<C^{\mathrm e}_v.
\]
For every integer \(n\ge2\),
\[
c^{\mathrm e}_v\rho_n^{\mathrm e}(v)
\le r_n^{*,\mathrm e}(v)
\le C^{\mathrm e}_v\rho_n^{\mathrm e}(v),
\qquad
B_n^{\mathrm e}\!\left(v,c^{\mathrm e}_v\rho_n^{\mathrm e}(v)\right)
\ge\frac25.
\]
The attaining rule satisfies
\[
\mathsf R_n^{\mathrm e}\!\left(P,\phi^{*,\mathrm e}_{n,v}\right)
\le0.01
\qquad(P\in H_0(v)).
\]
Whenever \(C^{\mathrm e}_v\rho_n^{\mathrm e}(v)<D_v\), it also satisfies
\[
1-\mathsf R_n^{\mathrm e}\!\left(P,\phi^{*,\mathrm e}_{n,v}\right)
\le0.01
\qquad
\left(P\in\mathcal M_v,\quad
C^{\mathrm e}_v\rho_n^{\mathrm e}(v)\le d(P)\right).
\]

For each such \(v,n\), there exist a null law \(P_0\in H_0(v)\) and a normalized finite prior
\[
\pi_1=\sum_{j=1}^{k}\omega_j\delta_{P_j},
\qquad
\omega_j\ge0,\qquad
\sum_{j=1}^{k}\omega_j=1,\qquad
\exists j:\ \omega_j>0,
\]
such that every positive-weight support law satisfies
\[
P_j\in\mathcal M_v,\qquad
e_{P_j}=e_{P_0},\qquad
c^{\mathrm e}_v\rho_n^{\mathrm e}(v)\le d(P_j)<D_v.
\]
Define its complete-record mixture by
\[
\mathbb P_{\pi_1,n}=\sum_{j=1}^{k}\omega_jP_j^{\otimes n},
\]
and let \(\operatorname{TV}\) denote the supremum of absolute event-probability differences. These laws satisfy
\[
\operatorname{TV}\!\left(P_0^{\otimes n},\mathbb P_{\pi_1,n}\right)<\frac12.
\]
For every supplied-propensity test \(\phi^{\mathrm e}\) satisfying
\[
\mathsf R_n^{\mathrm e}(P,\phi^{\mathrm e})\le\frac1{10}
\qquad(P\in H_0(v)),
\]
there is an index \(j\) with
\[
\omega_j>0,\qquad
1-\mathsf R_n^{\mathrm e}(P_j,\phi^{\mathrm e})\ge\frac25.
\]
For every integer \(n\ge N^{\mathrm e}_v\),
\[
C^{\mathrm e}_v\rho_n^{\mathrm e}(v)\le\frac{d_0}{2}.
\]

For \(P\in\mathcal M_v\), let \(Q_{a,P}(dy\mid x)\) denote the conditional outcome probability kernel in treatment arm \(a\). The untruncated inverse-propensity score identifies the original mean effect: for \(\lambda\)-almost every \(x\),
\[
e_P(x)\int_{\mathbb R}\frac{y}{e_P(x)}\,Q_{1,P}(dy\mid x)
+
\bigl(1-e_P(x)\bigr)
\int_{\mathbb R}\frac{-y}{1-e_P(x)}\,Q_{0,P}(dy\mid x)
=\tau_P(x).
\]
The frontier constants and nonempty-power thresholds are locally uniform over admissible tuples: on every compact subset of \(\mathcal V\), the lower multipliers admit a common positive lower bound, the upper multipliers admit a common finite upper bound, and the thresholds admit a common finite upper bound, including at \(p=2\) and \(\gamma=1/4\).
\end{theoremv}

\begin{theoremv}{thm:exact-propensity-preservation-boundary}[Exact propensity preservation boundary]
Let \(\mathcal V\) be the domain of public tuples \(v=(p,\alpha,\beta,\gamma)\) satisfying:
\begin{itemize}
\item \textbf{(Moment order.)} \(1<p\le2\).
\item \textbf{(Propensity smoothness.)} \(0<\alpha\le1\).
\item \textbf{(Baseline smoothness.)} \(0<\beta\le1\).
\item \textbf{(Effect smoothness.)} \(1/4\le\gamma\le1\).
\end{itemize}
For every \(v\in\mathcal V\), use the quantities and original-model radius multipliers \(c_v,C_v\) of \cref{obj:def:sharp-frontier-scales}, with
\[
q=\frac{p-1}{p},\qquad S=\alpha+\beta,\qquad
D_p=1+2\alpha+\frac{\beta}{q}+\frac{S}{2\gamma},
\]
\[
E_0(p)=\frac{2\gamma q}{2\gamma+q},\qquad
E_4(p)=\frac{2S}{D_p},\qquad
F(p)=\frac{2\alpha}{p-1}+\frac{\beta}{q}+\frac{S}{2\gamma}.
\]
Define the comparison exponent and the supplied-propensity radius multipliers by
\[
G(v)=\max\{0,E_0(p)-E_4(p)\},\qquad
c^{\mathrm e}_v=\frac{4^{-\gamma}}{16\sqrt3},\qquad
C^{\mathrm e}_v=3\bigl(120+128\sqrt{160\sqrt2}\bigr).
\]
Then
\[
G(v)=
\frac{2\gamma q}{(2\gamma+q)D_p}\max\{0,1-F(p)\}.
\]
For the unknown- and supplied-propensity experiments on the same class \(\mathcal M_v\) of \cref{obj:def:model}, the capped critical radii of \cref{obj:def:critical-radius,obj:def:oracle-radius} satisfy, for every integer \(n\ge2\),
\[
\frac{c_v}{C^{\mathrm e}_v}n^{G(v)}
\le
\frac{r_n^*(v)}{r_n^{*,\mathrm e}(v)}
\le
\frac{C_v}{c^{\mathrm e}_v}n^{G(v)}.
\]
If \(F(p)=1\), then
\[
0<\frac{c_v}{C^{\mathrm e}_v},
\qquad
0<\frac{C_v}{c^{\mathrm e}_v},
\]
and, for every integer \(n\ge2\),
\[
\frac{c_v}{C^{\mathrm e}_v}
\le
\frac{r_n^*(v)}{r_n^{*,\mathrm e}(v)}
\le
\frac{C_v}{c^{\mathrm e}_v}.
\]
If \(F(p)<1\), then
\[
\lim_{n\to\infty}\frac{r_n^*(v)}{r_n^{*,\mathrm e}(v)}=+\infty,
\]
and, for every integer \(n\ge2\),
\[
\frac{c_v}{C^{\mathrm e}_v}n^{E_0(p)-E_4(p)}
\le
\frac{r_n^*(v)}{r_n^{*,\mathrm e}(v)}
\le
\frac{C_v}{c^{\mathrm e}_v}n^{E_0(p)-E_4(p)}.
\]
The preservation region of \cref{obj:def:preservation-region} is exactly
\[
\mathcal R=
\left\{v\in\mathcal V:
\frac{2\alpha}{p-1}+\frac{\beta}{q}
+\frac{\alpha+\beta}{2\gamma}\ge1\right\}.
\]
Moreover, for every compact subset \(K\) of \(\mathcal V\), in the product topology on the four real coordinates, there exist real constants \(c,C\), with \(c>0\), such that, for every \(v\in K\),
\[
c\le\frac{c_v}{C^{\mathrm e}_v},
\qquad
64\sqrt3\,C_v\le C.
\]
\end{theoremv}
